% !TeX program = pdflatex
% Build: pdflatex, run until the cross-references converge (two passes are
% not enough, three are needed).
%
% English edition of version 3 (revision of 24 September 2026). It mirrors
% paper_ru_v3.tex section by section: the same \label keys, equations,
% numerical values and \cite keys, in the same order. Like paper_en.tex, it
% carries no GOST R 7.0.7-2021 apparatus (article type, UDC index, bilingual
% captions, second reference list, "Information about the author" block);
% instead it follows the layout expected by international journals:
% single-language captions, a Declarations block, and one reference list with
% DOIs. All numbers are those of paper_ru_v3.tex (analysis_v2.py with a fixed
% seed, tags [S0]-[S12]; stopping_rules.py, lpw_second_order.py and
% frontiers.py, tags [S13]-[S18], output in results_v3.txt).
\documentclass[11pt,a4paper]{article}

% cmap MUST precede fontenc: it embeds ToUnicode tables in the PDF, without
% which the ff/fi/fl ligatures cannot be extracted -- the text is then neither
% searchable nor correctly indexed by Crossref, Google Scholar and the like.
\usepackage{cmap}
\usepackage[T1]{fontenc}
\usepackage[utf8]{inputenc}
\usepackage[english]{babel}
\usepackage{amsmath,amssymb,amsthm}
\usepackage{booktabs}
\usepackage{graphicx}
% Figures are generated by analysis_v2.py; the English-labelled set lives in
% figs_v2/en/, the Russian one in figs_v2/ (figs/ holds the version-1 figures).
\graphicspath{{figs_v2/en/}}
\usepackage[margin=25mm]{geometry}
% xurl BEFORE hyperref: it allows long URLs to break anywhere. Without it,
% links such as t5k.org/mersenne/heuristic.html run into the right margin.
\usepackage{xurl}
\usepackage[colorlinks=true,allcolors=blue]{hyperref}
\emergencystretch=3em
\usepackage{xcolor}
\usepackage[expansion=false]{microtype}
\usepackage{caption}
\captionsetup{labelsep=period}

% Marker for quantities that cannot be obtained from the triples
% (N_b, n_min, n_max) and need the full index lists. All such markers were
% closed in version 3; the command is kept so that a stray marker still
% compiles (and shows up in red).
\newcommand{\recompute}[1]{\textcolor{red}{[recompute: #1]}}

\theoremstyle{plain}
\newtheorem{proposition}{Proposition}
\newtheorem{corollary}{Corollary}
\theoremstyle{remark}
\newtheorem{remark}{Remark}

\newcommand{\Rb}{R_b}
\newcommand{\eg}{e^{\gamma}}
\newcommand{\emg}{e^{-\gamma}}
\DeclareMathOperator{\Gam}{Gamma}
\DeclareMathOperator{\Pois}{Poisson}

\title{The observation scheme in the statistics of generalized\\
repunit primes: an exact conditional test of the\\
Lenstra--Pomerance--Wagstaff constant}
\author{Timofei S. Dokuchaev\thanks{Independent researcher, Irkutsk, Russia.
E-mail: \texttt{timofei93@yandex.ru}. ORCID:
\url{https://orcid.org/0009-0006-0510-5225}.}}
% Fixed date: \today would change on every rebuild.
\date{27 August 2026}

\begin{document}
\maketitle

\begin{abstract}
\noindent The Lenstra--Pomerance--Wagstaff (LPW) heuristic predicts
$\#\{n\le x: (b^n-1)/(b-1)\ \text{prime}\}\sim(e^{\gamma}/\ln b)\ln x$ and has been
tested since 1993 through a ``normalized occurrence frequency''. Right-censoring
does not bias this frequency: the endpoint form and the OLS slope are unbiased
exactly or up to an exponentially small term (Proposition~\ref{prop:unbiased}).
Instead, forms of the estimator disagree by up to $29.5\%$.

The main result concerns pooling $B=15$ bases ($M=218$ events). The search
frontiers are not recorded; they are Neyman--Scott nuisance parameters, and the
profile estimator $M/S$ is inconsistent. The position $D_b$ of the last find is
sufficient for the frontier, and conditionally on it
$N_b-1\sim\Pois(\lambda_bD_b)$ exactly (Proposition~\ref{prop:cond}). The exact
conditional likelihood gives the estimator $(M-B)/S$ instead of $(M-1)/S$ and
closed-form intervals: $\hat\kappa=1.838$, exact $95\%$ interval $[1.59,\,2.11]$,
$p=0.67$ against $e^{\gamma}\approx1.781$. The reported ten-percent excess is an
artifact of the scheme; the old test rejects a true hypothesis in $11$--$17\%$ of
cases.

This rests on the frontier being chosen without regard to the finds; otherwise
only an upper bound survives, and four documented frontiers show no violation in
the dangerous direction. The heuristic itself gives a finite-$n$ intensity
$(\kappa/\ln b)(1+c_b/t)$; $c_b$ is not determined by the data but moves
$\hat\kappa$ by $8$--$16\%$, and with it LPW fits the data and the ``slow
convergence from below'' is explained. The homogeneity test ($p=0.935$) detects
only a spread of $\kappa_b$ above $\pm37\%$.
\end{abstract}

\medskip
\noindent\textbf{Keywords:} repunit primes, generalized repunits, Mersenne
primes, Lenstra--Pomerance--Wagstaff conjecture, Euler--Mascheroni constant,
observation scheme, Type-I censoring, Poisson point process,
conditional likelihood, nuisance parameters, Neyman--Scott problem, sufficient
statistic, estimator bias, statistical power, computational number theory

\medskip
\noindent\textbf{Mathematics Subject Classification (2020):}
Primary 11Y11; Secondary 11N05, 62F03, 62B05, 62N01, 60G55, 11N80.

\section{Introduction}

For an integer base $b\ge2$ the generalized repunit is
\begin{equation}\label{eq:Rb}
\Rb(n)=\frac{b^{n}-1}{b-1},
\end{equation}
and $\Rb(n)$ can be prime only when $n$ is prime. For $b=2$ these are the
Mersenne numbers $M_p=2^{p}-1$. Building on the heuristics of
Wagstaff~\cite{wagstaff1983} and Pomerance~\cite{pomerance1981}, the LPW
conjecture asserts that
\begin{equation}\label{eq:LPW}
\pi_b(x)=\#\{n\le x:\ \Rb(n)\ \text{prime}\}\ \sim\ \frac{\eg}{\ln b}\,\ln x,
\end{equation}
where $\gamma$ is the Euler--Mascheroni constant; a detailed derivation of the
heuristic for $b=2$ is given in the PrimePages reference~\cite{t5k}. The base
enters~\eqref{eq:LPW} only through the factor $1/\ln b$: the constant
$\eg\approx1.7811$ is declared universal. Dubner~\cite{dubner1993} proposed
testing this universality empirically across bases, and
Bourdelais~\cite{bourdelais2009} has since maintained running estimates of the
normalized frequency, publishing them in the comments to OEIS
entries~\cite{oeis}. Related questions for the Mersenne sequence alone were
discussed recently in~\cite{wagstaff2025}. Close in spirit is
also~\cite{dominguez2026}, which looks for secondary structure in the Mersenne
exponents related to the factorization of $p-1$; methodologically it is built
like our~\S\ref{sec:struct}---stratified conditional likelihood plus a
significance test---but it tests for structure \emph{within} a single base,
whereas we test for it \emph{between} bases.

An interpretive folklore has grown up around these estimates. Base~13 is
described as ``lucky'' and base~7 as ``the unluckiest''; it is noted that the
normalized frequencies lie below the conjectured limit $\emg\approx0.5615$ and
converge to it slowly from below. The present work examines these claims, but
arrives at conclusions other than those one would naturally expect.

\subsection{Contributions}

\begin{enumerate}
\item\textbf{The frequency estimator in common use is not biased.} The scheme
in which the fitting interval ends at the largest known index does indeed
condition the estimate on the last observation being an event, and the
sufficient statistic is distributed as $\Gam(N,\lambda_b)$
(Proposition~\ref{prop:gamma}). But this does \emph{not} imply that the
normalized frequency is biased: both the endpoint form $D/(N\ln b)$ and the OLS
slope are exactly unbiased under scheme~A, and unbiased up to the factor
$1-(1+\mu)e^{-\mu}$ under scheme~B (Proposition~\ref{prop:unbiased},
Table~\ref{tab:mc}); at the observed $\mu\approx N_b\ge7$ the deficit does not
exceed $0.73\%$ for the endpoint form and $2.3\%$ for OLS without the origin.
Only the reciprocal $\hat\kappa=1/\hat G$ is biased, by the factor $N/(N-1)$,
and that is Jensen's inequality, not censoring. Unbiased estimators of $\kappa$
and of $1/\kappa$ cannot be reciprocals of one another; any ``correction'' that
repairs one of them damages the other (\S\ref{sec:notbias}).

\item\textbf{The genuine defect is instability with respect to the form of the
estimator.} Several numerically distinct variants are in circulation. On
identical data they disagree by as much as $29.5\%$ (base~18)---$1.7$ times the
median deviation of the estimates from $\emg$ ($17.6\%$), which is what gets
interpreted substantively; for two of the fifteen bases the choice of form
weighs more than their entire deviation from the LPW prediction. Under the
endpoint form base~18 is the third ``unluckiest'' of fifteen, and under OLS it
is the tenth, below the median. The label describes the choice of estimator, not
a property of the primes (\S\ref{sec:instability}).

\item\textbf{Main result: an exact conditional likelihood for scheme~B.} The
pivot $\kappa S\sim\Gam(M,1)$ is exact under the scheme ``observe until the
$N$-th event'', where $N$ is fixed by design. The data, however, arise from a
fixed frontier: the search covered $(t_0,\ln L_b]$ and the number of finds is
random. The unrecorded window widths $W_b$ are nuisance parameters whose number
grows with the number of bases; their joint MLE $\hat W_b=D_b$ leads to the
inconsistent profile estimator $M/S$---precisely the Neyman--Scott
situation~\cite{neymanscott1948}. We show (Proposition~\ref{prop:cond}) that
$D_b$ is sufficient for $W_b$ when $\kappa$ is held fixed and that,
conditionally on $D_b$, the quantity $N_b-1$ is $\Pois(\lambda_bD_b)$, with
\emph{the window width $W_b$ not entering this distribution}. Hence
$M-B\mid D\sim\Pois(\kappa S)$ is an exact Poisson pivot, and $(M-B)/S$ is the
conditional maximum-likelihood estimator in the sense of
Andersen~\cite{andersen1970}, exactly conditionally unbiased; it is also
obtained by integrating over $W_b$ against the measure $dW_b$
(Remark~\ref{rem:integrated}). No moment approximations, no exponentially small
remainders and no simulation-based inversion of intervals are required
(\S\ref{sec:scheme}). The distinction between the schemes is itself classical,
and the unbiasedness of $(N-1)/D$ under scheme~A is a textbook fact; we lay no
claim to either. What is new is that the frontier is \emph{undocumented}, so the
exposure must be reconstructed from the finds themselves, and the correction
turns out to be proportional to the number of bases
(Remark~\ref{rem:priority}).

\item\textbf{Numerical consequence.} The exact conditional inference gives
$\hat\kappa=1.838$ with a $95\%$ interval $[1.594,\,2.109]$ and $p=0.672$: the
excess over LPW disappears. Under scheme~B the old Gamma pivot rejects a true
hypothesis in $11$--$17\%$ of cases (\S\ref{sec:calib}); the spread depends on
how far the frontier has advanced, which is precisely why it is reported as a
range.

\item\textbf{The single assumption and its price.} The whole inference rests on
condition (A1): the frontier was chosen without regard to the finds. The
estimators $(M-1)/S$ and $(M-B)/S$ are built on the same statistics $(N_b,D_b)$
and differ by $6.9\%$; which of them is right is determined not by the data but
by how the data were collected. Without (A1) only an \emph{upper} bound holds
rigorously: by Wald's identity $\mathbb E[M]\ge\kappa\,\mathbb E[S]$ under any
stopping rule. There is no lower bound: the rule ``stop the search after a long
stretch without finds'' biases even $(M-B)/S$ upwards
(Remark~\ref{rem:bracket}, Table~\ref{tab:stop}). An attempt to sidestep the
choice by assembling the estimate from per-base pieces that are unbiased by
construction fails (\S\ref{sec:noescape}). Four documented frontiers
($b=2,10,22,26$) provide the first empirical check of (A1). The frontiers turn
out to lie closer to the last finds than (A1) predicts ($p=0.10$); that is, a
deviation, if there is one, points in the safe direction $m<1$
(\S\ref{sec:frontiers}). The only radical remedy is to publish the frontiers
$L_b$.

\item\textbf{Lucky bases: what the test cannot do.} Conditioning on $K=M-B$
makes the null distribution of the homogeneity test exact (multinomial), with
no appeal to $\chi^2$ asymptotics. We obtain $\Lambda=7.13$ on $14$ degrees of
freedom, exact $p=0.935$. But its power is such that it resolves only a spread
of $\kappa_b$ of order $\pm37\%$ (\S\ref{sec:homog}). The correct conclusion is
``base-dependent frequencies are not required by the data'', not ``they do not
exist''.

\item\textbf{The second-order term of the heuristic itself.} At finite $n$ the
LPW heuristic predicts not a constant intensity but
$\lambda_b(t)=(\kappa/\ln b)(1+c_b/t)$ with $c_b\in[\ln2,\,1.74]$: the divisors
of $\Rb(p)$ are at least $2p+1$. The exact conditional inference carries over
to this model in closed form (\S\ref{sec:second}). The data do not determine
$c_b$, but the estimate of $\kappa$ depends on it by $8$--$16\%$, that is, more
strongly than on the observation scheme. For $c=\ln2$ and $n>10$, LPW with its
own second-order term is consistent with the data ($p=0.63$--$0.74$), and the
trend of the estimate under truncation disappears. The same term explains the
folklore ``slow convergence from below''.

\item\textbf{An explicit power calculation.} With a two-sided test the pooled
sample resolves deviations from $\eg$ of at least $21.3\%$; the Mersenne
sequence alone, at least $45.5\%$ (\S\ref{sec:power}). The test convention is
stated explicitly, since the one-sided variant gives a noticeably different
threshold ($18.7\%$).
\end{enumerate}

Section~\ref{sec:sameerror} works through a narrower but instructive episode: a
naive regression of the Mersenne counting function produces an apparently
significant negative intercept, which turns out to be an artifact of
conditioning. We include it as a worked example of the same error in a different
guise.

\section{Model and data}\label{sec:data}

\subsection{The Poisson model}
Put $t=\ln n$. The heuristic~\eqref{eq:LPW} is equivalent to the statement that
the indices of generalized repunit primes to base~$b$ form a homogeneous Poisson
point process on the $t$-axis with constant intensity
\begin{equation}\label{eq:lam}
\lambda_b=\frac{\kappa}{\ln b},\qquad \text{LPW:}\ \ \kappa=\eg=1.781072\ldots
\end{equation}
Universality is the assertion that a single parameter $\kappa$ governs all
bases. The Poisson idealization here is a working assumption, not a theorem:
even for the ordinary primes the Poisson law in short intervals has been
established only conditionally, under the Hardy--Littlewood conjectures, and it
is known to break down as the scale grows~\cite{jha2026}. We adopt~\eqref{eq:lam}
as a model and treat $\kappa$ throughout as a quantity to be estimated; this
reduces the question ``is LPW true?'' to a one-parameter one and makes the bases
poolable.

\subsection{Admissible bases}\label{sec:bases}
Bases that are perfect powers are excluded: for $b=c^{m}$ with $m>1$ the number
$\Rb(n)$ admits an algebraic factorization and the process does not have the
assumed form. Bases $4,8,9,16$ are therefore omitted. Within the declared range
the selection does not depend on the results: we take \emph{all} bases
$2\le b\le20$ that are not perfect powers, which gives $B=15$; no base was
excluded or added on the strength of its estimate. Each of them has at least
eight known indices (the minimum is attained at $b=18$: eight terms of A133857).

We stress, however, that the upper limit $b\le20$ is itself a conventional round
number and not a consequence of the criterion ``at least eight indices'': that
criterion continues to hold beyond it (A127996 for $b=21$ has eight terms,
A127997 for $b=22$ has nine) and first fails at $b=23$ (A204940, six terms). To
make sure the limit does not look chosen after the fact, we tested its influence
directly. Extending the set to all bases $b\le26$ that are not perfect powers
(this adds $b=21,22,23,24,26$; $b=25=5^2$ is excluded algebraically) gives
$B=20$, $M=253$ and
\[
  \frac{M-1}{S}=1.9776,\qquad \frac{M-B}{S}=1.8285,
\]
that is, it shifts the main estimate~\eqref{eq:resB} by $0.5\%$---more than an
order of magnitude less than the price of the choice of observation scheme
(\S\ref{sec:scheme}). The homogeneity test on the extended set does not reject
homogeneity ($\Lambda=10.24$ on $19$ degrees of freedom, exact $p=0.953$;
tag~[S17]). No conclusion of this work depends on the limit $b\le20$; we retain
it as pre-declared.

\begin{remark}[Selection on the number of finds]\label{rem:selection}
Formally the limit $b\le20$ was declared in advance, but it was chosen in a
world where the numbers of known terms are already published, and this deserves
to be checked not rhetorically but through the sign of the shift. The
conditional inference of \S\ref{sec:scheme} conditions on $N_b\ge1$---which is
harmless, since every base has finds; conditioning on $N_b\ge7$, by contrast,
would not be harmless and would bias $\hat\kappa$ \emph{upwards}. Lifting the
limit to $b\le26$ adds, in particular, the base $b=23$ with six terms, and shifts
the estimate \emph{downwards}, from $1.8381$ to $1.8285$. The sign of the shift
is the one predicted by selection on the number of finds, and its size is an
order of magnitude smaller than the price of the observation scheme. We regard
this as sufficient evidence that the limit does not drive the conclusion.
\end{remark}

\subsection{The data and their status}
Table~\ref{tab:data} lists the sequences used, retrieved on 27 August 2026
(cross-check: Appendix~\ref{app:data}). Four caveats matter.

\paragraph{Probable primes.} For large indices most of the entries are probable
primes established by a Fermat or strong Fermat test after trial division. The
count is biased upwards only if a Fermat pseudoprime has been mistaken for a
prime, which for numbers of this size is practically unheard of. Sensitivity is
examined in~\S\ref{sec:robust}.

\paragraph{Verification frontiers.} For $b=2$ the GIMPS project~\cite{gimps}
publishes two frontiers: all exponents below $X_{\mathrm{ver}}=81\,648\,221$
have been tested twice, which gives exactly $50$ Mersenne primes, whereas all
exponents below $X_{\mathrm{ft}}=141\,308\,443$ have been tested at least once,
which gives $52$. The two largest exponents are provisional: unverified
candidates remain between them. The analysis of~\S\ref{sec:sameerror}, where the
frontier is used explicitly, is restricted to the $50$ verified exponents. The
pooled analysis is insensitive to this choice (\S\ref{sec:robust}).

Both frontiers are moving quantities, so we fix the moment of reading: the
PrimeNet Milestones report, 27 August 2026. No conclusion depends on the exact
value of $X_{\mathrm{ver}}$: it suffices that it lies between
$M_{50}=77\,232\,917$ and $M_{51}=82\,589\,933$, and this holds for every value
the report has taken during 2026.

\paragraph{The status of entry A000043.} At the retrieval date (27 Aug.\ 2026)
the DATA section of that entry contained $50$ terms and stopped at
$77\,232\,917$: OEIS regards the order of the subsequent exponents as unproven
until double-checking is complete. The exponents $82\,589\,933$ and
$136\,279\,841$ are known and documented by GIMPS (and are mentioned in the
comments to A000043), but formally they were not part of DATA. After
double-checking below $82\,589\,933$ was completed (GIMPS report of 4 Sep.\ 2026),
that exponent was added to DATA as the $51$st term (checked 24 Sep.\ 2026); this
does not affect the computations, since it is already included. We include both
exponents, because the statistic~\eqref{eq:pivotB} uses only the position
of the last event and not its ordinal number; \S\ref{sec:robust} verifies that
restricting base~2 to the fifty double-checked exponents does not change the
conclusion.

\paragraph{Search bounds and completeness of coverage for $b>2$.} For the other
bases a dated and attributed frontier above the last find has been published
only for $b=10$: ``Search limit is 10800000, currently'' (S.~Batalov, comment
on A004023, 1 July 2021). In the extended set it is joined by $b=22$ and $b=26$,
with the bound $10^5$ (R.~Price, comments on A127997 and A127999, 2012). For the
remaining bases the OEIS entries contain only the dates of the finds, and
participants in the distributed search state explicitly that the search limits
for most bases are unknown. It is precisely this gap that is the subject
of~\S\ref{sec:scheme} and, we believe, the principal obstacle to a quantitative
test of LPW; the four documented frontiers are used in~\S\ref{sec:frontiers}. We
note separately that the entire construction assumes that the search covered
$(t_0,\ln L_b]$ \emph{contiguously}. If the sieved ranges contain gaps---a common
occurrence in distributed searching---then $M$ is understated relative to the
exposure and $\hat\kappa$ is biased downwards, that is, in the same direction as
our main conclusion. In general, completeness of coverage cannot be checked from
published data. That such gaps are real is shown by base~$5$. In 2024 the
probable primes $R_5(4\,939\,471)$ and $R_5(5\,154\,509)$ were found
(R.~Propper); they are not in the OEIS, and summary tables mark the range
$3\,300\,593$--$4\,939\,471$ as not searched contiguously. For the conditional
inference this matters: law~\eqref{eq:condlaw} requires contiguous coverage of
$(t_0,t_0+D_b)$. We therefore retain $n_{\max}=3\,300\,593$. Including both
numbers, which is admissible only under contiguous coverage, would give
$\hat\kappa=1.852$ instead of $1.838$ (tag~[S18]). This is one more argument for
publishing not only the frontiers but also the sieving protocols.

\begin{table}[t]
\centering
\caption{Sequences used. $N_b$ is the number of events strictly to the right of
$t_0=\ln n_{\min}$, that is, one less than the number of known prime indices: the
smallest index sets the origin, not an event; $K_b=N_b-1$ is the conditional event
count entering~\eqref{eq:pivotB}. The $D_b$ column is rounded, but all computations
use exact $D_b=\ln n_{\max}-\ln n_{\min}$. In the last row $M=\sum_b N_b$, $K=M-B$
and $S=\sum_b D_b/\ln b$; the quantity $S$ is not the sum of the $D_b$ column.
Source: OEIS entries, retrieved 27.08.2026.}
\label{tab:data}
\small
\begin{tabular}{rlrrrrrrr}
\toprule
$b$ & OEIS & primes & $N_b$ & $K_b$ & $n_{\min}$ & $n_{\max}$ & $D_b$ & $D_b/\ln b$\\
\midrule
2 & A000043 & 52 & 51 & 50 & 2 & 136\,279\,841 & 18.0371 & 26.0220\\
3 & A028491 & 23 & 22 & 21 & 3 & 8\,530\,117 & 14.8605 & 13.5266\\
5 & A004061 & 19 & 18 & 17 & 3 & 3\,300\,593 & 13.9110 & 8.6434\\
6 & A004062 & 17 & 16 & 15 & 2 & 3\,360\,347 & 14.3344 & 8.0002\\
7 & A004063 & 10 & 9 & 8 & 5 & 1\,264\,699 & 12.4409 & 6.3934\\
10 & A004023 & 11 & 10 & 9 & 2 & 8\,177\,207 & 15.2237 & 6.6116\\
11 & A005808 & 13 & 12 & 11 & 17 & 1\,868\,983 & 11.6077 & 4.8408\\
12 & A004064 & 14 & 13 & 12 & 2 & 769\,543 & 12.8604 & 5.1754\\
13 & A016054 & 13 & 12 & 11 & 5 & 1\,503\,503 & 12.6139 & 4.9178\\
14 & A006032 & 11 & 10 & 9 & 3 & 1\,724\,417 & 13.2618 & 5.0252\\
15 & A006033 & 10 & 9 & 8 & 3 & 639\,833 & 12.2704 & 4.5311\\
17 & A006034 & 12 & 11 & 10 & 3 & 1\,990\,523 & 13.4053 & 4.7315\\
18 & A133857 & 8 & 7 & 6 & 2 & 1\,270\,141 & 13.3615 & 4.6228\\
19 & A006035 & 11 & 10 & 9 & 19 & 209\,359 & 9.3074 & 3.1610\\
20 & A127995 & 9 & 8 & 7 & 3 & 984\,349 & 12.7011 & 4.2397\\
\midrule
\multicolumn{3}{l}{total} & $M=218$ & $K=203$ & & & & $S=110.4423$\\
\bottomrule
\end{tabular}
\end{table}

\section{What is biased and what is not}\label{sec:notbias}

\subsection{The estimators in common use}
Dubner~\cite{dubner1993} and Bourdelais~\cite{bourdelais2009} characterize a
base~$b$ by a normalized occurrence frequency obtained by fitting a straight
line to $\ln n_k$ against $k$ and dividing by $\ln b$. By~\eqref{eq:lam} the
expected slope is $1/\lambda_b=\ln b/\kappa$, so the normalized frequency
estimates
\begin{equation}\label{eq:G}
G_b=\frac{1}{\kappa},\qquad\text{LPW:}\ \ G_b=\emg=0.5614595\ldots
\end{equation}
Several variants are in circulation---ordinary OLS with an intercept, regression
through the origin, and the endpoint form $\hat G_b=D_b/(N_b\ln b)$ with
$D_b=t_{N_b}-t_0$. All of them share one structural feature: the fitting
interval ends at the largest known index.

\subsection{The exact distribution}

\begin{proposition}[Conditional law of the right endpoint]\label{prop:gamma}
Let the events form a homogeneous Poisson process of intensity $\lambda$ on
$(t_0,\infty)$, and let $t_1<\dots<t_N$ be its first $N$ events after $t_0$.
Then $D:=t_N-t_0$ has the $\Gam(N,\lambda)$ distribution with density
$\lambda^{N}d^{N-1}e^{-\lambda d}/(N-1)!$. Conditionally on $D=d$, the preceding
$N-1$ events are distributed as the order statistics of $N-1$ independent
uniform variables on $(t_0,t_0+d)$.
\end{proposition}

\begin{proof}
The waiting times $t_1-t_0,\dots,t_N-t_{N-1}$ are independent
$\mathrm{Exp}(\lambda)$ variables, and a sum of $N$ of them is
$\Gam(N,\lambda)$. Conditional uniformity is the order-statistics property of
the Poisson process, applied to $(t_0,t_0+d)$.
\end{proof}

\begin{corollary}[Bias of the reciprocal]\label{cor:kappa}
The maximum-likelihood estimator of $\lambda$ based on $D$ alone is
$\hat\lambda=N/D$, and
\begin{equation}\label{eq:bias}
\mathbb E[\hat\lambda]=N\,\mathbb E[D^{-1}]=\frac{N}{N-1}\,\lambda,
\end{equation}
since $\mathbb E[D^{-1}]=\lambda/(N-1)$ for $D\sim\Gam(N,\lambda)$ with $N\ge2$.
The estimator $\tilde\lambda=(N-1)/D$ is exactly unbiased.
\end{corollary}

\subsection{What does not follow from Corollary~\ref{cor:kappa}}

It is natural to conclude that, since $\hat\lambda$ is inflated by the factor
$N/(N-1)$, the frequency $\hat G_b$ must be deflated by the same factor. This is
false under either of the two observation schemes analysed below.

\begin{proposition}[Unbiasedness of the frequency]\label{prop:unbiased}
\emph{(a) Scheme~A.} Under the hypotheses of Proposition~\ref{prop:gamma}:
\begin{enumerate}
\item[\rm(i)] the endpoint estimator is exactly unbiased:
$\mathbb E\!\left[\dfrac{D}{N\ln b}\right]=\dfrac{1}{\lambda\ln b}=G_b$;
\item[\rm(ii)] the OLS slope in the regression of $t_k$ on $k$ is exactly
unbiased: $\mathbb E[\hat\beta]=1/\lambda$;
\item[\rm(iii)] consequently the ``corrected'' quantity
$\tilde G_b=\frac{N}{N-1}\hat G_b$ is biased \emph{upwards}, by exactly the
factor $N/(N-1)$.
\end{enumerate}
\emph{(b) Scheme~B.} Let the process be observed on a window of width $W$, and
put $\mu=\lambda W$; the estimators are set to zero where they are undefined.
Then
\begin{equation}\label{eq:BunbiasedB}
\mathbb E\!\left[\frac{D}{N}\,\mathbf 1_{\{N\ge1\}}\right]
=\mathbb E\bigl[\hat\beta_0\,\mathbf 1_{\{N\ge1\}}\bigr]
=\frac{1}{\lambda}\Bigl(1-(1+\mu)e^{-\mu}\Bigr),
\end{equation}
where $\hat\beta_0$ is the OLS slope through the points $(k,t_k)$,
$k=0,\dots,N$, \emph{including} the origin. For the slope $\hat\beta_1$ through
the points $k=1,\dots,N$ \emph{without} the origin, which is defined for
$N\ge2$,
\begin{equation}\label{eq:BunbiasedOLS}
\mathbb E\bigl[\hat\beta_1\bigm|N\ge2\bigr]
=\frac{1}{\lambda}\cdot
\frac{1-\bigl(1+\mu+\mu^2/2\bigr)e^{-\mu}}{1-(1+\mu)e^{-\mu}} .
\end{equation}
Both forms are thus unbiased up to a term exponentially small in $\mu$; at
$\mu=7$ the deficit is $0.73\%$ in~\eqref{eq:BunbiasedB} and $2.25\%$
in~\eqref{eq:BunbiasedOLS}, and at $\mu=10$ it is $0.05\%$ and $0.23\%$.
\end{proposition}

\begin{proof}
(a)(i) $\mathbb E[D]=N/\lambda$ for $D\sim\Gam(N,\lambda)$, whence the result at
once. (ii) The OLS slope is the linear statistic $\hat\beta=\sum_k w_k t_k$ with
weights $w_k=(k-\bar k)/\sum_j(j-\bar k)^2$, for which $\sum_k w_k=0$ and
$\sum_k w_k k=1$. Since $\mathbb E[t_k]=t_0+k/\lambda$, we get
$\mathbb E[\hat\beta]=\sum_k w_k\,(t_0+k/\lambda)=1/\lambda$.
(iii) is immediate from (i).

(b) Under scheme~B, conditionally on $N=n\ge1$, the positions of the events are
the order statistics of $n$ independent uniform variables on $(t_0,t_0+W]$, so
$D/W\sim\mathrm{Beta}(n,1)$ and $\mathbb E[D\mid N=n]=Wn/(n+1)$, whence
$\mathbb E[D/N\mid N=n]=W/(n+1)$. Averaging over $N\sim\mathrm{Pois}(\mu)$ and
taking into account that $D$ is defined only for $N\ge1$,
\[
\mathbb E\Bigl[\frac{D}{N}\Bigr]
=W\sum_{n\ge1}\frac{\mu^{n}e^{-\mu}}{n!\,(n+1)}
=\frac{W}{\mu}\sum_{n\ge1}\frac{\mu^{n+1}e^{-\mu}}{(n+1)!}
=\frac{W}{\mu}\Bigl(1-e^{-\mu}-\mu e^{-\mu}\Bigr),
\]
which is~\eqref{eq:BunbiasedB} for the endpoint form. For the OLS slope, note
that $\mathbb E[t_k\mid N=n]=t_0+kW/(n+1)$ for $k=0,\dots,n$; since the slope is
a linear statistic with weights $\sum_kw_k=0$, $\sum_kw_kk=1$, in both variants
$\mathbb E[\hat\beta\mid N=n]=W/(n+1)$. For $\hat\beta_0$ the summation runs over
$n\ge1$ and gives the same result as above. For $\hat\beta_1$ it runs over
$n\ge2$: $\sum_{n\ge2}\frac{W}{n+1}\frac{\mu^ne^{-\mu}}{n!}
=\frac{W}{\mu}\bigl(1-(1+\mu+\mu^2/2)e^{-\mu}\bigr)$, and division by
$\Pr(N\ge2)=1-(1+\mu)e^{-\mu}$ gives~\eqref{eq:BunbiasedOLS}.
\end{proof}

The substantive point is simple: unbiased estimators of $\kappa$ and of
$1/\kappa$ cannot be reciprocals of one another, because $x\mapsto1/x$ is
strictly convex. The unbiased pair is $(\hat G_b,\ \tilde\kappa_b)$ with
$\tilde\kappa_b=(N_b-1)\ln b/D_b$; but $\tilde\kappa_b\ne1/\hat G_b$. Any single
``censoring correction'' that repairs one of these quantities damages the other.
Table~\ref{tab:mc} confirms Proposition~\ref{prop:unbiased} by simulation under
both schemes.

\begin{remark}\label{rem:schemeindep}
It is essential that part~(b) removes a possible objection to the first
conclusion of this work. Below (\S\ref{sec:scheme}) we adopt scheme~B as
matching the way the data were obtained; if part~(a) were the only one, the
whole of \S\ref{sec:notbias} would rest on a scheme that the paper itself
rejects. Part~(b) shows that the conclusion does not depend on the choice of
scheme: the deficit $(1+\mu)e^{-\mu}$ at $\mu\approx N_b$ is $0.73\%$ for
$N_b=7$, $0.05\%$ for $N_b=10$ and about $2\cdot10^{-6}$ for $N_b=16$ (for OLS
without the origin,~\eqref{eq:BunbiasedOLS}: $2.25\%$, $0.23\%$ and
$1.4\cdot10^{-5}$), whereas the deviations discussed in the literature are tens
of percent.
\end{remark}

\begin{table}[t]
\centering
\caption{Monte Carlo check of Proposition~\ref{prop:unbiased}, $60\,000$
replications per row (seed fixed, see \texttt{analysis\_v2.py}, tag~[S1]). The
target value of every column except the second is $1.0000$. In the right-hand
part $\mu=\lambda W$; the column ``theory'' gives the factor
$1-(1+\mu)e^{-\mu}$ of~\eqref{eq:BunbiasedB} for $\hat G$. The scheme~B column
$\mathbb E[\hat\beta]\lambda$ is for the slope $\hat\beta_1$ without the origin,
conditional on $N\ge2$; its theoretical value is given
by~\eqref{eq:BunbiasedOLS}: $0.9775$, $0.9977$ and $1.0000$ at $\mu=7,10,16$.}
\label{tab:mc}
\small
\begin{tabular}{rrrrrr@{\qquad}rrr}
\toprule
& \multicolumn{5}{c}{scheme A ($N$ fixed)} & \multicolumn{3}{c}{scheme B ($\mu=\lambda W$)}\\
\cmidrule(r){2-6}\cmidrule(l){7-9}
$N=\mu$ & $\mathbb E[\hat\lambda]/\lambda$ & $N/(N-1)$
& $\mathbb E[\tilde\lambda]/\lambda$ & $\mathbb E[\hat G]\lambda$
& $\mathbb E[\hat\beta]\lambda$
& $\mathbb E[\hat G]\lambda$ & theory & $\mathbb E[\hat\beta]\lambda$\\
\midrule
7 & 1.1655 & 1.1667 & 0.9990 & 1.0012 & 1.0020 & 0.9908 & 0.9927 & 0.9790\\
10 & 1.1103 & 1.1111 & 0.9992 & 1.0007 & 1.0018 & 0.9997 & 0.9995 & 0.9948\\
16 & 1.0650 & 1.0667 & 0.9985 & 1.0016 & 1.0017 & 0.9992 & 1.0000 & 0.9992\\
22 & 1.0462 & 1.0476 & 0.9987 & 1.0011 & 1.0006 & 0.9991 & 1.0000 & 0.9992\\
51 & 1.0206 & 1.0200 & 1.0006 & 0.9994 & 0.9992 & 0.9997 & 1.0000 & 0.9999\\
\bottomrule
\end{tabular}
\end{table}

\begin{remark}
The negative result of this subsection deserves emphasis, because it runs
against intuition: the systematic downward deviation of the published
frequencies from $\emg$ is \emph{not} explained by the fitting scheme. Its
source is not the arithmetic of the estimator but the second-order term of the
heuristic itself, $c_b\ln\ln x$ in the counting function, which depresses the
frequency at finite $n$ and vanishes only as $1/\ln n$ (\S\ref{sec:second}); the
observation scheme also contributes to the gap (\S\ref{sec:scheme}).
\end{remark}

\subsection{Instability with respect to the form of the estimator}\label{sec:instability}

Unbiasedness does not mean agreement. Table~\ref{tab:bases} gives the endpoint
form $\hat G_b$ and the OLS estimate $\hat G_b^{\mathrm{OLS}}$ on identical
data. The discrepancies reach $29.5\%$ ($b=18$), $15.9\%$ ($b=13$) and $10.3\%$
($b=11$).

Compare these with what the literature interprets substantively---the deviations
of the estimates themselves from $\emg$. Across the fifteen bases the latter
have median $17.6\%$ (from $5.6\%$ at $b=20$ to $43.7\%$ at $b=19$). The largest
discrepancy between forms, $29.5\%$, exceeds that median by a factor of $1.7$.
On average the choice of form is of course cheaper: the median discrepancy is
$4.1\%$, about a quarter of the median deviation. But for two of the fifteen
bases ($b=18$ and $b=20$) the choice of form weighs more than their entire
deviation from the LPW prediction---and one of them, base~18, is precisely a base
the folklore discusses.

The consequence for the folklore is direct. Under the endpoint form base~18 has
$\hat G_{18}=0.6604$, the third largest of fifteen (after $b=7$ with $0.7104$
and $b=10$ with $0.6612$), so it belongs firmly among the three ``unluckiest'';
under OLS, $\hat G_{18}^{\mathrm{OLS}}=0.4656$, tenth of fifteen and below the
median $0.4904$, so it has become ``lucky''. The label flips without a single
new prime, purely from a change in the form of the estimator: the base moves
from third place to tenth.

To this is added a second, technical difficulty with OLS: its standard errors
are wrong in this problem. The residuals are deviations of order statistics and
are strongly autocorrelated; for $b=2$ we obtain $\rho_1=0.825$, and the
Newey--West covariance estimator with eight lags raises the standard error of
the intercept from $0.447$ to $0.732$ (\S\ref{sec:sameerror}). Published point
values come with no intervals at all, and this---not bias---is the principal
methodological objection to current practice.

\subsection{The conditional likelihood under scheme A}\label{sec:condlik}

Proposition~\ref{prop:gamma} yields more than an analysis of bias. Setting
$t_0=\ln n_{\min}$ and taking $(N_b,D_b)$ as the data for base~$b$, we obtain
the log-likelihood
\begin{equation}\label{eq:loglik}
\ell_{\mathrm A}(\kappa)=\sum_b\Bigl[N_b\ln\frac{\kappa}{\ln b}-\frac{\kappa}{\ln b}D_b\Bigr]+\mathrm{const},
\end{equation}
whose distribution does not depend on the search bound $L_b$. This matters in
practice: $L_b$ is documented only for $b=2$ and $b=10$. Differentiating
\eqref{eq:loglik} and putting $S=\sum_b D_b/\ln b$, we get
\begin{equation}\label{eq:mle}
\hat\kappa=\frac{M}{S},\qquad \tilde\kappa=\frac{M-1}{S},\qquad M=\sum_b N_b .
\end{equation}
Since $\kappa S=\sum_b\lambda_bD_b$ is a sum of independent $\Gam(N_b,1)$
variables,
\begin{equation}\label{eq:pivot}
\kappa S\sim\Gam(M,1),
\end{equation}
an exact pivot; confidence intervals and $p$-values follow in closed form.
\emph{It is precisely this statement that requires revision}---see
\S\ref{sec:scheme}.

\begin{remark}[Independence across bases]\label{rem:indep}
Pooling presupposes that the processes for different bases are independent.
Strictly this fails at small indices: seventeen index values occur in more than
one base, and $n=3$ occurs in nine of them, since $\Rb(3)=b^2+b+1$ is prime for
many~$b$. All such coincidences lie at $n\le317$. More important than coinciding
indices (the numbers $\Rb(n)$ for different $b$ are distinct) is the possible
sharing of algebraic divisors; both affect only small~$n$, whereas $D_b$ is
determined by the upper end of the range. Section~\ref{sec:robust} confirms that
discarding all indices below~$10^3$ leaves the conclusion unchanged. We note that
this remark bears more strongly on the conditional multinomial test of
\S\ref{sec:homog} than on the pooled estimate: dependence distorts not the mean
but the null distribution of the statistic; the same truncation check covers
that test as well.
\end{remark}

\section{Two observation schemes}\label{sec:scheme}

\subsection{The distinction}

The pivot~\eqref{eq:pivot} was derived on the assumption that $N_b$ is fixed by
the observation plan: we watch the process until $N_b$ events have accumulated,
then stop. Call this \emph{scheme~A} (inverse sampling, Type-II censoring).

The data arose differently. For each base the computational search covered the
whole interval $(t_0,\ln L_b]$, and every event falling inside it was found; the
number of finds $N_b$ is a random variable, while the frontier $L_b$ is fixed.
Call this \emph{scheme~B}.

\begin{remark}[What here is classical and what is not]\label{rem:priority}
The distinction between the schemes is itself classical and goes back to
reliability theory, where scheme~A is called a \emph{failure-truncated} and
scheme~B a \emph{time-truncated} life test~\cite{epstein1954trunc,%
epsteinsobel1954}. The estimator $(N-1)/D$, unbiased under scheme~A
(Corollary~\ref{cor:kappa}), is a textbook fact there, and we lay no claim to
it. The distribution of the MLE of the mean under Type-I censoring was studied
by Bartholomew~\cite{bartholomew1963}. We also note an elementary fact: under
scheme~B with a known window the sufficient statistic is
$N\sim\Pois(\lambda W)$, and $\mathbb E f(N)=\sum_nf(n)e^{-\mu}\mu^n/n!$ is an
entire function of $\mu=\lambda W$, so it cannot coincide with $W/\mu$, which has
a pole at zero. Hence no unbiased estimator of the \emph{mean} $1/\lambda$ exists
under scheme~B (for time-censored samples from the exponential distribution
see~\cite{sengupta2009}); the same argument applies conditionally on $D_b$.
Proposition~\ref{prop:cond} does not contradict this: it gives exact
unbiasedness of an estimator of the \emph{intensity}.

What we regard as new is something else. In the classical setting the moment at
which the test is stopped is \emph{known}: the time $T$ is set by the
experimenter. Here the frontier $L_b$ is documented only for $b=2$ and $b=10$
(\S\ref{sec:frontiers}), so the full exposure cannot be read off a protocol and
must instead be reconstructed---via the memorylessness of the process---from the
finds themselves. That is where a correction \emph{per base}, rather than one for
the sample as a whole, comes from: there are as many unknown frontiers as there
are bases. A survey of the biases generated by the edge of the observation
window, though in a nonparametric setting and for the shape of the interevent
distribution rather than for the intensity, is given in~\cite{kivela2015}.

Dual problems have long been known. In paleontology the unknown end of a
stratigraphic range is estimated from fossil finds, and the observed extreme
horizon is extended by roughly one mean gap between
finds~\cite{straussSadler1989,marshall1990}. In ecology the unknown extinction
time of a species is estimated from its sighting record~\cite{solow1993}. The
general problem of estimating a truncation point is treated
in~\cite{robsonwhitlock1964}. There, however, the unknown end of the window is
the parameter of interest. Here it is a nuisance parameter, and it is the
intensity that is estimated. Closest in technique is~\cite{liu2023}: there the
observation window starts at the first event, and the likelihood factorizes into
a Poisson factor and a factor that does not depend on the intensity. We have
found no work in which Type-I censoring with an unrecorded censoring time is
treated as a problem of pooled estimation of an intensity across many strata.
\end{remark}

The distinction is not cosmetic. Under scheme~B, conditionally on $N_b=n$, the
positions of the events are $n$ independent uniform variables on
$(t_0,\ln L_b]$, so the conditional distribution of $D_b$ depends on $L_b$ and
\emph{not} on $\lambda_b$: with $N_b$ fixed, $D_b$ is informative about the
frontier, not about the intensity.

More precisely, under scheme~B with a known window width $W_b=\ln L_b-t_0$ the
likelihood for base~$b$ is $\lambda_b^{N_b}e^{-\lambda_bW_b}$, whereas with an
unknown width it depends on the nuisance parameter $W_b\ge D_b$ and decreases
monotonically in it. The joint MLE therefore sets $\hat W_b=D_b$---the frontier
exactly at the last find---and expression~\eqref{eq:loglik} turns out to be the
\emph{profile} likelihood over the nuisance parameters $W_1,\dots,W_B$, with
$M/S$ the profile MLE. The number of nuisance parameters grows with the number
of bases, and the profile MLE is inconsistent: as $B\to\infty$ with
$\lambda_bW_b$ bounded, it does not converge to $\kappa$. This is the classical
Neyman--Scott situation~\cite{neymanscott1948} (for a survey
see~\cite{lancaster2000}), and the standard way out is to condition on a
statistic sufficient for the nuisance
parameters~\cite{kalbfleischsprott1970,andersen1970}. This is what
\S\ref{sec:exact} implements.

Note also that \S\ref{sec:sameerror} of the present work, where the frontier for
$b=2$ is known and used explicitly, rests on scheme~B. One cannot apply both
schemes in a single paper; below we adopt scheme~B as matching the way the data
were obtained, and construct exact inference for it.

\subsection{The exact conditional law under scheme B}\label{sec:exact}

The key observation is that under scheme~B one should condition not on the
number of events but on the position of the last of them. The unknown frontier
then disappears from the problem entirely. One assumption is required.

\medskip
\noindent\textbf{(A1)} \emph{The position of the frontier $L_b$ was not chosen in
the light of the search results: the computational budget, not the finds,
determines how far in $n$ the sieving has been carried. Formally: the window
width $W_b=\ln L_b-t_0$ is a random variable with an arbitrary unknown
distribution $\pi_b$ (in particular, a constant) that is independent of the
process of finds for base~$b$.}

\medskip
\noindent In the language of survival analysis, (A1) is the requirement that the
censoring be \emph{non-informative}: the stopping time does not depend on what
the process has produced. The distributions $\pi_b$ may be arbitrary and may
differ between bases; they are never estimated.

\begin{proposition}[Exact conditional likelihood]\label{prop:cond}
Suppose that under scheme~B a homogeneous Poisson process of intensity
$\lambda_b$ is observed on $(t_0,t_0+W_b]$, that (A1) holds and $N_b\ge1$, and
let $D_b=t_{N_b}-t_0$ be the distance from the origin to the last event found.
\begin{enumerate}
\item[\rm(i)] The joint density of $(N_b,D_b)$ at $(n,d)$, $n\ge1$, $d>0$,
factorizes:
\begin{equation}\label{eq:factor}
p(n,d)=\underbrace{e^{-\lambda_bd}\frac{(\lambda_bd)^{n-1}}{(n-1)!}}_{\text{depends on }\lambda_b}
\;\cdot\;
\underbrace{\lambda_b\int_{[d,\infty)}e^{-\lambda_b(w-d)}\,\pi_b(dw)}_{\text{density of }D_b},
\end{equation}
so that for fixed $\lambda_b$ the statistic $D_b$ is sufficient for $W_b$ (and
for the whole distribution $\pi_b$).
\item[\rm(ii)] Consequently,
\begin{equation}\label{eq:condlaw}
K_b:=N_b-1\ \bigm|\ D_b\ \sim\ \Pois(\lambda_b D_b),
\end{equation}
and this distribution depends neither on $W_b$ nor on $\pi_b$.
\item[\rm(iii)] For processes that are independent across bases,
\begin{equation}\label{eq:pivotB}
K:=M-B\ \bigm|\ (D_b)_{b}\ \sim\ \Pois(\kappa S),\qquad
S=\sum_b\frac{D_b}{\ln b},
\end{equation}
and $\hat\kappa_{\mathrm B}=(M-B)/S$ is the conditional maximum-likelihood
estimator, exactly conditionally unbiased:
$\mathbb E[\hat\kappa_{\mathrm B}\mid (D_b)_b]=\kappa$.
\end{enumerate}
\end{proposition}

\begin{proof}
(i) First fix $W_b=w$. For $0<d<d+\Delta\le w$ and $\Delta\to0$, the event
$\{N_b=n,\ D_b\in[d,d+\Delta)\}$ means: exactly $n-1$ events in $(t_0,t_0+d)$,
one in $[t_0+d,t_0+d+\Delta)$ and none in $[t_0+d+\Delta,t_0+w]$. By the
independence of the increments of the Poisson process its probability equals
$e^{-\lambda_bd}\frac{(\lambda_bd)^{n-1}}{(n-1)!}\cdot\lambda_b\Delta\cdot
e^{-\lambda_b(w-d)}+o(\Delta)$, and for $d>w$ it is zero. By (A1), averaging
over $w\sim\pi_b$ affects only the last factor and gives~\eqref{eq:factor}.
Summation over $n\ge1$ shows that the second factor in~\eqref{eq:factor} is the
marginal density of $D_b$.
(ii) Divide~\eqref{eq:factor} by the marginal density of $D_b$.
(iii) A sum of independent Poisson variables is Poisson; the conditional
log-likelihood $\ell_{\mathrm B}(\kappa)=\sum_b\bigl[K_b\ln\kappa-\kappa D_b/\ln b\bigr]
+\mathrm{const}$ is maximized at $\kappa=K/S$, and conditional unbiasedness is
immediate from $\mathbb E[K\mid D]=\kappa S$.
\end{proof}

\begin{remark}[Conditional and integrated likelihood]\label{rem:integrated}
The same result is obtained by integrating the likelihood over the nuisance
parameter. For a fixed window the joint density is proportional to
$\lambda_b^{n}d^{\,n-1}e^{-\lambda_bW_b}\mathbf 1_{\{W_b\ge d\}}$, and
$\int_d^\infty\lambda_b^{n}e^{-\lambda_bW_b}\,dW_b=\lambda_b^{n-1}e^{-\lambda_bd}$,
which is exactly the kernel of~\eqref{eq:condlaw}. This is the integrated
likelihood in the sense of~\cite{berger1999} with an improper uniform measure on
$W_b$. Such a procedure depends, in itself, on the choice of measure: a measure
uniform in $\mu_b=\lambda_bW_b$, that is, $d\mu_b=\lambda_b\,dW_b$, would give
$\lambda_b^{n}e^{-\lambda_bd}$ and bring back the inconsistent estimator $M/S$.
Here, however, the choice is not arbitrary. First, the result coincides with the
\emph{conditional} likelihood~\eqref{eq:condlaw}, which is defined without any
measure and is invariant under reparametrization of $W_b$. Second, $W_b$ enters
\eqref{eq:factor} as a location parameter of the quantity $W_b-D_b$, and $dW_b$
is the invariant measure of the translation group. Third, for $\pi_b(dw)=dw$ the
marginal factor in~\eqref{eq:factor} equals
$\lambda_b\int_d^\infty e^{-\lambda_b(w-d)}dw=1$ and does not depend on
$\lambda_b$. Hence, when the frontier is unknown, the factor discarded by
conditioning carries no information about $\kappa$ that could be used without
specifying $\pi_b$.
\end{remark}

\begin{remark}[Connection with the Neyman--Scott problem]\label{rem:neymanscott}
The correction ``one event per base'' is of the same nature as the loss of one
degree of freedom per group in the classical Neyman--Scott
example~\cite{neymanscott1948}. For $y_{ij}\sim N(\mu_i,\sigma^2)$ the profile
MLE $\hat\sigma^2$ is inconsistent, whereas conditioning on $\bar y_i$, which is
sufficient for $\mu_i$ at fixed $\sigma^2$, gives
$\sum_j(y_{ij}-\bar y_i)^2\sim\sigma^2\chi^2_{n_i-1}$. In both cases the nuisance
parameter of a group is one-dimensional and is a location parameter, and the
statistic sufficient for it ``absorbs'' exactly one observation: the mean absorbs
one degree of freedom, the last event absorbs one event. The general mechanism is
the conditional likelihood given a statistic that is sufficient for the nuisance
parameter at a fixed value of the parameter of
interest~\cite{kalbfleischsprott1970}. Its consistency as the number of strata
grows was established under regularity conditions by
Andersen~\cite{andersen1970}, and the optimality of the conditional estimator in
models that mix over the nuisance parameter by Lindsay~\cite{lindsay1980}.
\end{remark}

\begin{remark}[Inhomogeneous intensity]\label{rem:inhom}
Argument~(i) does not use homogeneity. If the intensity $\lambda_b(t;\theta)$
depends on a parameter $\theta$, then conditionally on $D_b$ the events in
$(t_0,t_0+D_b)$ form a Poisson process of the same intensity, and the
conditional log-likelihood is
\begin{equation}\label{eq:condinhom}
\ell_{\mathrm B}(\theta)=\sum_b\Bigl[\sum_{i=1}^{N_b-1}\ln\lambda_b(t_i;\theta)
-\int_{t_0}^{t_0+D_b}\lambda_b(t;\theta)\,dt\Bigr]:
\end{equation}
the last event is excluded from the sum, and the integral stops at it. It is in
this form that the variable-intensity models of \S\ref{sec:struct} should be
estimated.
\end{remark}

The practical consequences of Proposition~\ref{prop:cond} deserve to be listed
separately, since they remove almost all the caveats of the previous version of
this work.

\begin{enumerate}
\item \emph{Exact intervals without simulation.} Inverting the Poisson rate
gives the Garwood interval~\cite{garwood1936}
\begin{equation}\label{eq:garwood}
\Bigl[\tfrac{1}{2S}\chi^2_{\alpha/2,\,2K},\ \tfrac{1}{2S}\chi^2_{1-\alpha/2,\,2(K+1)}\Bigr],
\end{equation}
whose coverage is never below nominal. For our $S$ it equals $0.9540$ at
$\kappa=\eg$ and lies within $0.950$--$0.958$ for $\kappa\in[1.5,\,2.2]$ (the
excess is the usual effect of discreteness).
\item \emph{No remainder.} No exponentially small terms, no caveats about bases
without finds: conditioning on $N_b\ge1$ holds by construction, since $D_b$ is
defined only then. A base without a single find does not enter the conditional
likelihood.
\item \emph{No dependence on the position of the frontier.} The rule by which
$L_b$ was chosen does not enter~\eqref{eq:condlaw} at all---provided (A1) holds.
\item \emph{The cost of conditioning is small.} $\operatorname{SD}(\hat\kappa_{\mathrm B})
=\sqrt{\kappa/S}=0.1270$ against $\sqrt{\kappa/S_{\mathrm{full}}}=0.1224$ with
known frontiers standing at their expected positions,
$S_{\mathrm{full}}\approx S+B/\kappa$ (Proposition~\ref{prop:delta}). Forgoing
knowledge of $L_b$ increases the standard deviation by $3.7\%$, or costs $7.1\%$ of
the Fisher information: $B/\kappa$ out of $S+B/\kappa$, that is, roughly $B$ events
out of $M$. Knowledge of the frontiers is needed not so much for precision as for
checking (A1) (\S\ref{sec:noescape}).
\end{enumerate}

\noindent If (A1) is violated, conditioning on $D_b$ ceases to be harmless. If
the search stops exactly at a find, then $D_b$ is itself the stopping time,
$\delta_b=0$, and we are back at scheme~A. If instead the search stops after a
sufficiently long stretch without finds, then the stopping itself carries
information about the process before $D_b$, and law~\eqref{eq:condlaw} no longer
holds (Remark~\ref{rem:bracket}).

\subsection{What ignorance of the frontier costs}\label{sec:cost}

Besides Remark~\ref{rem:neymanscott}, the size of the correction is explained by
a direct moment calculation. We retain it as a substantive interpretation and as
the basis for the analysis of violations of (A1) in Remark~\ref{rem:bracket}.

\begin{proposition}\label{prop:delta}
Suppose that under scheme~B the search for base~$b$ covered $(t_0,\ln L_b]$,
that the window width $W_b=\ln L_b-t_0$ is fixed (under (A1), conditionally on
$W_b$), and let $\delta_b=\ln L_b-t_{N_b}$ be the distance from the frontier to
the last event found (with $\delta_b=W_b$ when $N_b=0$). Then
\begin{equation}\label{eq:delta}
\Pr(\delta_b>x)=e^{-\lambda_b x}\ \ (0\le x<W_b),\qquad
\Pr(\delta_b=W_b)=e^{-\lambda_b W_b},
\end{equation}
that is, $\delta_b$ is an exponential variable with parameter $\lambda_b$
truncated at $W_b$, and
\begin{equation}\label{eq:Edelta}
\mathbb E[\delta_b]=\frac{1-e^{-\lambda_b W_b}}{\lambda_b}\ \le\ \frac{1}{\lambda_b}.
\end{equation}
The full exposure is
$S_{\mathrm{full}}=S+\sum_b\delta_b/\ln b=\sum_b W_b/\ln b$, with
$\mathbb E[M]=\kappa S_{\mathrm{full}}$ identically, and
\begin{equation}\label{eq:EBd}
\mathbb E\Bigl[\kappa\sum_b\delta_b/\ln b\Bigr]=B-\sum_b e^{-\lambda_b W_b}.
\end{equation}
\end{proposition}

\begin{proof}
For $x<W_b$ the event $\{\delta_b>x\}$ means that the interval
$(\ln L_b-x,\ln L_b]$ of length $x$ contains no event of the process, which has
probability $e^{-\lambda_b x}$. As $x\to W_b^-$ what remains is the case ``no
events at all'', which gives the atom in~\eqref{eq:delta}. Integrating the tail,
$\mathbb E[\delta_b]=\int_0^{W_b}e^{-\lambda_b x}\,dx=(1-e^{-\lambda_b W_b})/\lambda_b$,
whence~\eqref{eq:Edelta}; the inequality is obvious. Next, $S_{\mathrm{full}}$ is
not random: by construction $D_b+\delta_b=W_b$, so
$\kappa S_{\mathrm{full}}=\sum_b\lambda_b W_b=\mathbb E[M]$---this is exactly the
moment equation. Finally
$\kappa\,\mathbb E[\delta_b]/\ln b=\lambda_b\mathbb E[\delta_b]=1-e^{-\lambda_b W_b}$,
and summation over $b$ gives~\eqref{eq:EBd}.
\end{proof}

In other words, ignorance of the frontier costs information equivalent to one
event \emph{per base}---whereas formula~\eqref{eq:mle} subtracts one event
\emph{in total}. With $B=15$ the difference is fourteen events, that is, $6.9\%$:
\begin{equation}\label{eq:cost}
\frac{M-1}{S}=1.9648\qquad\text{versus}\qquad\frac{M-B}{S}=1.8381 .
\end{equation}
The remainder $\sum_b e^{-\lambda_b W_b}$ dropped in~\eqref{eq:EBd} is the
expected number of bases for which the search produced no find at all; in our
data $\lambda_bW_b\approx N_b\ge7$, so it does not exceed $3\cdot10^{-3}$
against $B=15$. We stress that in the exact formulation~\eqref{eq:condlaw} this
remainder does not arise at all: the moment calculation is needed here only for
interpretation.

\begin{remark}[What remains without (A1)]\label{rem:bracket}
The previous version claimed that the family $(M-Bm)/S$ with $m\in[0,1]$ yields a
provable two-sided bracket $[(M-B)/S,\ (M-1)/S]$. This claim is \emph{false},
and we withdraw it. Inequality~\eqref{eq:Edelta} was derived for a fixed window
(or one independent of the process), that is, precisely under (A1), and then no
bracket is needed: the estimator $(M-B)/S$ is correct. If (A1) is violated,
$\lambda_b\mathbb E[\delta_b]$ can be any non-negative number.

Only the \emph{upper} bound survives rigorously. If the frontier is a stopping
time with respect to the search history (the decision to stop does not look
ahead), then by Wald's identity for the Poisson process
$\mathbb E[N_b]=\lambda_b\,\mathbb E[W_b]$ \emph{under any} stopping rule with
$\mathbb E[W_b]<\infty$. Since $W_b=D_b+\delta_b\ge D_b$, we get
$\mathbb E[M]\ge\kappa\,\mathbb E[S]$: in the sense of moments $M/S$ bounds
$\kappa$ from above. There is no assumption-free lower bound. The rule ``stop
the search after a stretch of $g$ mean intervals without finds'' gives
$\lambda_b\delta_b\ge g$. Moreover, waiting for such a stretch itself packs the
events before $D_b$ more densely, so that law~\eqref{eq:condlaw} no longer holds.
Table~\ref{tab:stop} shows the consequences: at $g=2$ the estimator $(M-B)/S$ is
inflated by $6.7\%$, and the exact test of Proposition~\ref{prop:cond} rejects a
true hypothesis in $18\%$ of cases instead of $5\%$. The rule ``stop at the next
find'' ($\delta_b=0$) biases $(M-B)/S$ downwards by $6.4\%$, and then the
estimator $(M-1)/S$ is the correct one.

The precise statement is therefore as follows. Under (A1) the main
result~\eqref{eq:resB} is exact. Under stopping ``at a find'' the
estimator~\eqref{eq:resA} is correct, and it also bounds $\kappa$ from above in
the sense of moments under any rule. Under stopping ``after a drought'' the true
$\kappa$ is \emph{below} $(M-B)/S$, that is, even closer to $\eg$ or below it.
None of the violations of (A1) creates an excess over LPW that the data would
require explaining. But without documented frontiers the size of the
uncertainty is not bounded from below: this is a stronger argument for
publishing $L_b$ than the one given previously.
\end{remark}

\begin{table}[t]
\centering
\caption{Scheme~B under various stopping rules: relative mean of the two pooled
estimators, the ratio $\mathbb E[M]/(\kappa\mathbb E[S])$ (at least $1$ under every
rule by Wald's identity), and the type~I error of the exact two-sided Poisson
test~\eqref{eq:pivotB} at the $5\%$ level. Budget $W^0_b=N_b/\lambda_b$;
``drought $g$'': after the budget is spent the search continues until
$g/\lambda_b$ passes without a find. $20\,000$ replications per row, accuracy
about $\pm0.05\%$ for the means and $\pm0.003$ for the level (script
\texttt{stopping\_rules.py}, tag~[S13], seed $20260924$).}
\label{tab:stop}
\small
\begin{tabular}{lrrrr}
\toprule
Stopping rule & $\mathbb E[(M{-}1)/S]/\kappa$ & $\mathbb E[(M{-}B)/S]/\kappa$
& $\dfrac{\mathbb E[M]}{\kappa\mathbb E[S]}$ & type~I error\\
\midrule
(A1): frontier $=W^0_b$ & 1.0688 & 0.9999 & 1.0737 & 0.046\\
(A1): frontier $=UW^0_b$, $U\sim\mathrm U(0.5,\,2)$ & 1.0551 & 1.0000 & 1.0582 & 0.043\\
stop at a find & 0.9963 & 0.9362 & 1.0003 & 0.165\\
drought $g=0.5$ & 1.0762 & 1.0074 & 1.0811 & 0.050\\
drought $g=1$ & 1.0923 & 1.0250 & 1.0973 & 0.063\\
drought $g=2$ & 1.1227 & 1.0667 & 1.1277 & 0.178\\
\bottomrule
\end{tabular}
\end{table}

\subsection{Calibration of the old pivot}\label{sec:calib}

How bad the Gamma pivot~\eqref{eq:pivot} is when the data are generated by
scheme~B depends on how far the frontier has advanced---and this should be
reported as a range rather than as a single number. Table~\ref{tab:calib} gives
both quantities at once: the type~I error of the old test and the relative bias
of both pooled estimators.

\begin{table}[t]
\centering
\caption{Scheme~B under seven frontier rules: type~I error of the nominal $5\%$
test built on the Gamma pivot~\eqref{eq:pivot}, and the relative bias of the two
pooled estimators. The rules are stated through the window width
$W_b=\ln L_b-t_0$: ``$\mathbb E[\#]$ = observed'' is $W_b=N_b/\lambda_b$;
``$+s$ mean intervals'' is $W_b=(N_b+s)/\lambda_b$; ``$c$ times further in $t$''
is $W_b=c\,D_b$ with the observed $D_b$. $50\,000$ replications per row, accuracy
about $\pm0.003$ for the level and $\pm0.1\%$ for the means (tag~[S10]).}
\label{tab:calib}
\small
\begin{tabular}{lrrr}
\toprule
Rule for the frontier & type~I error & $\mathbb E[(M{-}1)/S]/\kappa$
& $\mathbb E[(M{-}B)/S]/\kappa$\\
\midrule
$\mathbb E[\#\text{events}]$ = observed & 0.1634 & 1.0691 & 1.0001\\
frontier $+0.5$ mean intervals & 0.1620 & 1.0667 & 1.0002\\
frontier $+1$ mean interval & 0.1596 & 1.0640 & 0.9998\\
frontier $+2$ mean intervals & 0.1527 & 1.0599 & 0.9998\\
frontier $+5$ mean intervals & 0.1359 & 1.0506 & 1.0002\\
frontier $1.5$ times further in $t$ & 0.1342 & 1.0499 & 0.9999\\
frontier twice as far in $t$ & 0.1136 & 1.0371 & 1.0001\\
\midrule
A (control) & 0.0492 & --- & ---\\
\bottomrule
\end{tabular}
\end{table}

Three conclusions follow from Table~\ref{tab:calib}. First, under scheme~B a
nominal $5\%$ test rejects a true hypothesis in $11$--$17\%$ of cases: the
claimed precision of the intervals is overstated two- to threefold, and the more
so the closer the frontier is to the last find. Second, the estimator $(M-1)/S$
is biased upwards by $4$--$7\%$ under all rules, whereas $(M-B)/S$ is unbiased
up to Monte Carlo noise under all rules---as Proposition~\ref{prop:cond}
predicts, since conditional unbiasedness implies unconditional unbiasedness.
Third, the figure ``$16.2\%$'' quoted in an earlier version is the value for one
particular scenario; the correct thing is to report the range. All seven rules
of Table~\ref{tab:calib} place the frontier without regard to the finds, that
is, they satisfy (A1); rules that violate (A1) are analysed in
Remark~\ref{rem:bracket} and Table~\ref{tab:stop}.

\subsection{Documented frontiers: an empirical check of (A1)}\label{sec:frontiers}

A dated and attributed frontier above the last find could be found for four
bases (checked on 24 September 2026):
\begin{itemize}
\item $b=2$: the GIMPS report at the time the data were read,
$X_{\mathrm{ver}}=81\,648\,221$ and $X_{\mathrm{ft}}=141\,308\,443$
(\S\ref{sec:data});
\item $b=10$: ``Search limit is 10800000, currently'' (S.~Batalov, A004023,
1 July 2021);
\item $b=22$ and $b=26$: ``No other terms < 100000'' (R.~Price, A127997,
A127999, 2012; the last finds, $27\,823$ and $26\,717$, were made in
2007--2008).
\end{itemize}
The other mentions of search limits are either outdated and lie far below
$n_{\max}$, or coincide with the last find and mark it rather than the search.

Under (A1) the quantities $\lambda_b\delta_b$ are independent and distributed as
$\mathrm{Exp}(1)$ (the truncation at $W_b$ is negligible here), so their sum over
$k$ bases has the $\Gam(k,1)$ distribution. Small values point to stopping ``at a
find'' ($m<1$), large ones to stopping ``after a drought'' ($m>1$,
Remark~\ref{rem:bracket}). The observed values are $0.143$ ($b=2$ with
$X_{\mathrm{ver}}$; $0.093$ with $X_{\mathrm{ft}}$), $0.215$ ($b=10$), $0.737$
($b=22$) and $0.722$ ($b=26$). Their mean is $0.44$--$0.45$ against an expected
value of one. The probability of so small a sum under (A1) is $0.10$--$0.11$
(tag~[S18]).

The conclusions must be cautious: there are four observations, the $b=2$
frontier moves continuously, and two of the bases enter only the extended set.
Nevertheless this is the first check of (A1) against data, and its result is
informative about the \emph{direction}. The documented frontiers lie not farther
from but closer to the last finds than (A1) predicts. The violation that is
dangerous for the main conclusion---stopping ``after a drought'', which biases
$(M-B)/S$ upwards---is not supported by the data: the upper tail is
$0.89$--$0.90$. If (A1) is violated at all, it is in the direction $m<1$, and the
true value lies between $(M-B)/S$ and $(M-1)/S$.

The documented frontiers can also be used directly. For bases with a known
frontier the full likelihood applies: $N_b\sim\Pois(\lambda_bW_b)$ with all
events and the full window. For the others the conditional
law~\eqref{eq:condlaw} applies. The sum is again exactly Poisson. The hybrid
estimate with the known frontiers for $b=2$ ($X_{\mathrm{ft}}$) and $b=10$ is
$\hat\kappa=1.853$ with an exact $95\%$ interval $[1.608,\,2.125]$ and
$p=0.588$ against $\eg$. It differs from $(M-B)/S=1.838$ by $0.8\%$, and the
main conclusion does not change.

\subsection{The data do not choose the scheme for us}\label{sec:noescape}

We stress that $(M-1)/S$ and $(M-B)/S$ are functions of the same sufficient
statistics $(N_b,D_b)$. No processing of the available data will say which of
them is right: the choice is determined by how the data were collected, and its
price is $6.9\%$---twice the entire residual gap with LPW ($3.2\%$).

The natural attempt to sidestep the choice is to build the pooled estimate out
of pieces that are unbiased by construction. The per-base estimator
$\tilde\kappa_b=(N_b-1)\ln b/D_b$ is unbiased under scheme~A
(Corollary~\ref{cor:kappa}) and unbiased under scheme~B up to an exponentially
small term. Indeed, for $N_b=n\ge2$ the ratio $D_b/W_b$ is distributed as
$\mathrm{Beta}(n,1)$, and the identity
$\mathbb E[1/\mathrm{Beta}(\alpha,\beta)]=(\alpha+\beta-1)/(\alpha-1)$ gives
$\mathbb E[\tilde\lambda_b\mid N_b=n]=n/W_b$; for $n\in\{0,1\}$ the estimator is
zero. Averaging over $n\sim\mathrm{Pois}(\lambda_bW_b)$ therefore gives not
$\lambda_b$ but
\[
  \mathbb E[\tilde\lambda_b]=\lambda_b-\frac{\Pr(N_b=1)}{W_b}
  =\lambda_b\Bigl(1-e^{-\lambda_bW_b}\Bigr),
\]
that is, the same exponentially small deficit: for $\lambda_bW_b\approx N_b\ge7$
it does not exceed $0.1\%$.

Nevertheless, the weighted mean
$$\kappa_\ast=\sum\nolimits_b(N_b-2)\tilde\kappa_b\Big/\sum\nolimits_b(N_b-2)$$
with inverse-variance weights has a bias of $0.0\%$ under scheme~A but
$+7.5\%$ under scheme~B, that is, practically the same as $(M-1)/S$. The reason
is that the weights---the event counts themselves---are random under scheme~B
and positively correlated with the estimates; unbiasedness of each summand does
not carry over to a ratio of sums. The correct way to pool is not to weight the
per-base estimates but to sum the sufficient statistics, as~\eqref{eq:pivotB}
does.

From this we draw a narrow but, we think, important conclusion: \emph{the gain
in precision from pooling the bases is inseparable from an assumption about the
observation scheme}. The only radical remedy is to document the search frontiers
$L_b$, as GIMPS does for $b=2$. We address this request to the participants of
the relevant search projects.

\section{Results}\label{sec:results}

\subsection{Per-base estimates}

\begin{table}[t]
\centering
\caption{Per-base estimates. $\hat G_b=D_b/(N_b\ln b)$ is the endpoint form;
$\hat G_b^{\mathrm{OLS}}$ is the OLS slope of $\ln n_k$ on $k$ divided by $\ln b$;
both are unbiased (Proposition~\ref{prop:unbiased}).
$\hat\kappa_b=K_b/(D_b/\ln b)$ is the scheme~B conditional estimator with the
exact Garwood interval~\eqref{eq:garwood} from $\Pois(\kappa D_b/\ln b)$. The last
column is the Poisson mid-PIT at $\kappa=\eg$: values near one correspond to bases
with a \emph{high} $\kappa_b$ (orientation opposite to the Gamma PIT of the earlier
version). LPW predicts $G=\emg=0.5615$ and $\kappa=\eg=1.7811$.}
\label{tab:bases}
\small
\begin{tabular}{rrrrrrrl r}
\toprule
$b$ & $N_b$ & $K_b$ & $\hat G_b$ & $\hat G_b^{\mathrm{OLS}}$ & diff., \% &
$\hat\kappa_b$ & 95\% CI (Garwood) & mid-PIT\\
\midrule
2 & 51 & 50 & 0.5102 & 0.5309 & $+4.1$ & 1.921 & [1.426, 2.533] & 0.710\\
3 & 22 & 21 & 0.6148 & 0.6081 & $-1.1$ & 1.552 & [0.961, 2.373] & 0.272\\
5 & 18 & 17 & 0.4802 & 0.4666 & $-2.8$ & 1.967 & [1.146, 3.149] & 0.670\\
6 & 16 & 15 & 0.5000 & 0.4904 & $-1.9$ & 1.875 & [1.049, 3.092] & 0.594\\
7 & 9 & 8 & 0.7104 & 0.7371 & $+3.8$ & 1.251 & [0.540, 2.466] & 0.160\\
10 & 10 & 9 & 0.6612 & 0.6746 & $+2.0$ & 1.361 & [0.622, 2.584] & 0.216\\
11 & 12 & 11 & 0.4034 & 0.3617 & $-10.3$ & 2.272 & [1.134, 4.066] & 0.794\\
12 & 13 & 12 & 0.3981 & 0.3796 & $-4.7$ & 2.319 & [1.198, 4.050] & 0.820\\
13 & 12 & 11 & 0.4098 & 0.3448 & $-15.9$ & 2.237 & [1.117, 4.002] & 0.780\\
14 & 10 & 9 & 0.5025 & 0.5386 & $+7.2$ & 1.791 & [0.819, 3.400] & 0.528\\
15 & 9 & 8 & 0.5035 & 0.5054 & $+0.4$ & 1.766 & [0.762, 3.479] & 0.513\\
17 & 11 & 10 & 0.4301 & 0.4353 & $+1.2$ & 2.114 & [1.014, 3.887] & 0.717\\
18 & 7 & 6 & 0.6604 & 0.4656 & $-29.5$ & 1.298 & [0.476, 2.825] & 0.228\\
19 & 10 & 9 & 0.3161 & 0.2971 & $-6.0$ & 2.847 & [1.302, 5.405] & 0.911\\
20 & 8 & 7 & 0.5300 & 0.5713 & $+7.8$ & 1.651 & [0.664, 3.402] & 0.444\\
\bottomrule
\end{tabular}
\end{table}

Eleven of the fifteen bases lie below $\emg$ under the endpoint form, and eleven
of the fifteen under OLS---but these are \emph{different} elevens. The membership
changes: base~18 lies above $\emg$ under the endpoint form and below it under
OLS, and base~20 the other way round. In other words, even the sign of the
deviation is not fully robust to the choice of form; only their \emph{number} is.
This is the same effect as in \S\ref{sec:instability}, and it affects exactly the
two bases for which the choice of form weighs more than the deviation itself.

All fifteen mid-PIT quantiles lie inside the central $95\%$ region; the extremes
are $b=19$ at quantile $0.911$ and $b=7$ at $0.160$. Base~13, the ``lucky'' one,
sits at $0.780$, which is in no way remarkable for one base out of fifteen. All
fifteen Garwood intervals cover $\eg$.

\subsection{The pooled estimate}\label{sec:pooled}

The exact conditional inference for scheme~B (Proposition~\ref{prop:cond})
gives
\begin{equation}\label{eq:resB}
\hat\kappa_{\mathrm B}=\frac{M-B}{S}=1.8381,\qquad
95\%\ \text{CI}=[1.5939,\,2.1091],\qquad p=0.672,
\end{equation}
where the interval is~\eqref{eq:garwood} with $K=203$, $S=110.4423$, and the
$p$-value is obtained by the doubled-minimum-tail convention for
$K\sim\Pois(\eg S)$, $\eg S=196.71$; under the mid-$p$ convention $p=0.647$.
The exact coverage of interval~\eqref{eq:resB} is $0.9540$ at $\kappa=\eg$ and
does not fall below $0.950$ for any~$\kappa$.

For comparison, scheme~A via~\eqref{eq:mle}--\eqref{eq:pivot} gives
\begin{equation}\label{eq:resA}
\tilde\kappa=1.9648,\qquad 95\%\ \text{CI}=[1.7205,\,2.2444],\qquad p=0.142 .
\end{equation}
The point estimate~\eqref{eq:resA} is $10.3\%$ above $\eg$, and it is precisely
this excess that used to be interpreted as the trace of a factor omitted by the
heuristic. On moving to scheme~B the excess falls to $3.2\%$, and the $p$-value
rises from $0.14$ to $0.67$.

We hold that scheme~B is closer to the way the data were obtained, and therefore
regard~\eqref{eq:resB} as the main result. Estimate~\eqref{eq:resA} corresponds
to stopping the search exactly at a find and, by Remark~\ref{rem:bracket}, bounds
$\kappa$ from above in the sense of moments under any stopping rule; under
stopping ``after a drought'' the true value lies below~\eqref{eq:resB}. Without
publication of the $L_b$, (A1) cannot be checked. Both estimates are consistent
with LPW; they differ in whether the data call for an explanation of a residual
excess. We conclude that they do not. The overall picture is shown in
Figure~\ref{fig:kappa}: the confidence interval of every base covers both $\eg$
and both pooled estimates.

\begin{figure}[t]
\centering
\includegraphics[width=.85\textwidth]{fig_kappa.pdf}
\caption{Per-base estimates $\hat\kappa_b$ with exact $95\%$ Garwood
intervals~\eqref{eq:garwood}; marker area is proportional to $N_b$. Dashed line: the
LPW value $\eg$; green with a band: the pooled scheme~B estimate and its interval;
blue dash-dotted: the scheme~A estimate $(M-1)/S$. Each base separately is
compatible with a universal constant.}
\label{fig:kappa}
\end{figure}

\subsection{Homogeneity, and its power}\label{sec:homog}

Let $H_0$ be the model with a single $\kappa$ and $H_1$ the model with free
$\kappa_b$. From the conditional likelihood of Proposition~\ref{prop:cond}, since
$\sum_b\hat\kappa_bD_b/\ln b=K=\hat\kappa_{\mathrm B}S$, the likelihood-ratio
statistic reduces to
\begin{equation}\label{eq:LR}
\Lambda=2\sum_b K_b \ln\frac{\hat\kappa_b}{\hat\kappa_{\mathrm B}}
=2\sum_b K_b\ln\frac{K_b}{K\,s_b/S},\qquad s_b=\frac{D_b}{\ln b}.
\end{equation}
Crucially, its null distribution is \emph{exact} here: conditioning on
$K=\sum_bK_b$ turns~\eqref{eq:pivotB} into the multinomial model
$(K_b)\mid K\sim\mathrm{Multinomial}(K;\,s_b/S)$, and there is no need to appeal
to $\chi^2_{B-1}$ asymptotics. We obtain $\Lambda=7.127$, exact $p=0.935$
($2\cdot10^6$ multinomial replications, $\pm0.0002$); the asymptotic
approximation gives $p=0.930$ on $14$ degrees of freedom, that is, it works well
in this case.

But a non-rejection is informative exactly to the extent that the test is
powerful, and here one must apply to oneself the same requirement one applies to
the literature. We simulated alternatives in which the $\kappa_b$ are lognormally
scattered about $\eg$ with coefficient of variation $cv$:

\begin{center}
\small
\begin{tabular}{lrrrrrrrrr}
\toprule
$cv$ & 0.05 & 0.10 & 0.15 & 0.20 & 0.25 & 0.30 & 0.35 & 0.40 & 0.50\\
power & 0.059 & 0.097 & 0.184 & 0.326 & 0.484 & 0.645 & 0.774 & 0.851 & 0.948\\
\bottomrule
\end{tabular}
\end{center}

The test attains $80\%$ power only at $cv\approx0.37$ ($12\,000$ replications
per grid point, power accurate to about $\pm0.004$). Hence $p=0.935$ means: the
data do not require base-dependent frequencies, but neither do they exclude a
spread of $\kappa_b$ exceeding a third of the value of the constant. The
statement ``lucky bases do not exist'' does not follow from the available data;
what follows is the weaker statement ``their existence is not supported by the
data, and differences below $\pm37\%$ are in principle indistinguishable''.

\subsection{Robustness}\label{sec:robust}

\paragraph{Lower truncation.} We varied the lower truncation point in order to
probe the asymptotic regime (the heuristic~\eqref{eq:LPW} is asymptotic, whereas
the constant-intensity model starts at $n=2$). The results are in
Table~\ref{tab:trunc}. The scheme~A estimate grows from $1.965$ to $2.099$---not
strictly monotonically (at $n>10^{2}$ it retreats from $2.019$ to $2.012$) but
systematically: the growth is $6.8\%$, and all five values come from nested
subsamples of the same data. This is a trend, not noise. The scheme~B estimate
shows no upward trend and stays close to $\eg$, within $1.789$--$1.871$ (under
the weaker dropout criterion $N_b\ge1$, within $1.757$--$1.871$), and its
$p$-values lie between $0.53$ and $1.00$ at all truncation levels. The reason is
visible from~\eqref{eq:pivotB}: truncation reduces $M$ but not $B$, whereas
scheme~A subtracts one regardless of the sample size. We regard this as
confirmation of the diagnosis of \S\ref{sec:scheme}: the upward trend was a
property of the estimator, not of the data. From $n>10$ onwards the scheme~B
estimate, on the contrary, \emph{decreases} monotonically ($1.871\to1.789$). This
is exactly what the second-order term of the heuristic itself requires, and once
that term is taken into account the decrease disappears (\S\ref{sec:second}).
Both curves are shown in Figure~\ref{fig:trunc}.

\begin{table}[t]
\centering
\caption{Robustness to lower truncation. A base is retained if at least
\emph{two} events survive the truncation ($N_b\ge2$); at $n>10^4$ base~19 no longer
meets this criterion, whence $B=14$. The criterion matters: under the weaker
requirement $N_b\ge1$ base~19 is retained, giving the last row. The interval and
$p_{\mathrm B}$ are exact, scheme~B; the last column ($p_{\mathrm A}$, the scheme~A
Gamma pivot) is given for comparison with the old practice.}
\label{tab:trunc}
\small
\begin{tabular}{lrrrrlrr}
\toprule
truncation & $B$ & $M$ & $(M-1)/S$ & $(M-B)/S$ & 95\% CI, scheme B
& $p_{\mathrm B}$ & $p_{\mathrm A}$\\
\midrule
none & 15 & 218 & 1.9648 & 1.8381 & [1.594, 2.109] & 0.672 & 0.142\\
$n>10$ & 15 & 192 & 2.0192 & 1.8712 & [1.606, 2.168] & 0.531 & 0.080\\
$n>10^2$ & 15 & 156 & 2.0116 & 1.8299 & [1.540, 2.158] & 0.771 & 0.124\\
$n>10^3$ & 15 & 125 & 2.0551 & 1.8231 & [1.498, 2.197] & 0.832 & 0.106\\
$n>10^4$ & 14 & 89 & 2.0990 & 1.7889 & [1.407, 2.242] & 1.000 & 0.116\\
\midrule
$n>10^4$, $N_b\ge1$ & 15 & 90 & 2.0848 & 1.7569 & [1.382, 2.202] & 0.966 & 0.128\\
\bottomrule
\end{tabular}
\end{table}

\begin{figure}[t]
\centering
\includegraphics[width=.75\textwidth]{fig_truncation.pdf}
\caption{The pooled estimate under lower truncation: $(M-1)/S$ (scheme~A),
$(M-B)/S$ (scheme~B) and $K/S(c)$ with the second-order term $c=\ln2$
(\S\ref{sec:second}). The upward trend of scheme~A disappears under scheme~B; the
residual decline of scheme~B disappears once the second-order term is included.}
\label{fig:trunc}
\end{figure}

\paragraph{Leave-one-base-out.} A leave-one-base-out analysis gives
$\hat\kappa_{\mathrm B}$ within $[1.808,\,1.878]$ (the extreme values arise when
$b=19$ and $b=3$ are removed) and exact $p$-values within
$[0.49,\,0.85]$; no single base determines the conclusion.

\paragraph{Status of the data.} Restricting base~2 to the fifty double-checked
exponents moves the result from $\tilde\kappa=1.965$, $p=0.142$ to
$\tilde\kappa=1.961$, $p=0.151$ (scheme~A) and leaves the scheme~B estimate
practically unchanged. Excluding all indices $n>10^6$, where the entries have
PRP status, gives $\tilde\kappa=1.951$ with $M=182$. In the worst case, when the
last known PRP of some base turns out to be composite, the largest shift comes
from $b=17$: $\tilde\kappa=1.977$. All variants lie inside
interval~\eqref{eq:resA}, and a fortiori inside~\eqref{eq:resB}, so possible
pseudoprimes do not affect the conclusion.

\section{Statistical power}\label{sec:power}

By~\eqref{eq:pivotB} the power against a true constant $\kappa=r\eg$ at level
$\alpha$ is available in closed form (Table~\ref{tab:power}). The convention has
to be stated: throughout this work the test is \emph{two-sided}, with the
doubled-minimum-tail convention, $\alpha=0.05$. The one-sided variant of the same
test gives noticeably different numbers; they are shown in the right-hand part
of the table for comparability with works in which the convention is not stated.

\begin{table}[t]
\centering
\caption{Power at level $\alpha=0.05$ for rejecting $\kappa=\eg$ when the true
constant exceeds it by the stated factor. Exposures: $S_2=26.02$ ($K=50$) for base~2
alone and $S=110.44$ ($K=203$) for the pooled sample.}
\label{tab:power}
\small
\begin{tabular}{lrr@{\qquad}rr}
\toprule
& \multicolumn{2}{c}{two-sided} & \multicolumn{2}{c}{one-sided}\\
\cmidrule(r){2-3}\cmidrule(l){4-5}
$\kappa/\eg$ & Mersenne only & pooled & Mersenne only & pooled\\
\midrule
1.05 & 0.060 & 0.099 & 0.082 & 0.165\\
1.10 & 0.099 & 0.266 & 0.147 & 0.386\\
1.15 & 0.164 & 0.514 & 0.235 & 0.644\\
1.20 & 0.253 & 0.752 & 0.343 & 0.844\\
1.30 & 0.479 & 0.972 & 0.581 & 0.988\\
1.50 & 0.861 & 1.000 & 0.910 & 1.000\\
\bottomrule
\end{tabular}
\end{table}

At $80\%$ power the Mersenne sequence alone detects only deviations of $45.5\%$
and above; pooling fifteen bases lowers that threshold to $21.3\%$ (one-sided
variant: to $18.7\%$). Attaining $80\%$ power against a $10\%$ deviation requires
an exposure $S\approx469$, that is, $K\approx835$ conditional events
($M\approx850$) at $\kappa=\eg$ (if the true constant is $10\%$ larger, there
will be $\approx919$ events); against $5\%$, $S\approx1816$ and $K\approx3234$.
The power function of a discrete test is not monotone in $S$ (it is
saw-toothed), so the values given are those beyond which the power no longer
drops below $80\%$.

These thresholds are computationally out of reach, and it is worth saying by how
much. Since $K$ grows logarithmically in the search frontier, doubling the
frontiers \emph{of all fifteen bases} yields $\eg\ln2\sum_b1/\ln b=9.64$
additional events for the entire pooled sample. To gather the missing $632$
events one needs $632/9.64\approx65.6$ doublings, that is, a growth of the
frontiers by a factor of $2^{65.6}\approx5.5\cdot10^{19}$; for a sensitivity of
$5\%$, about $314$ doublings, that is, a factor of order $10^{94}$. This is not a
matter of decades but of magnitudes lying beyond the reach of any conceivable
computation. Both power curves, as functions of the true deviation, are shown in
Figure~\ref{fig:power}.

\begin{figure}[t]
\centering
\includegraphics[width=.72\textwidth]{fig_power.pdf}
\caption{Power of the two-sided rejection of $\kappa=\eg$ at level $\alpha=0.05$.
The horizontal line marks $80\%$ power.}
\label{fig:power}
\end{figure}

Hence a reading rule for this literature. A report that the data agree with LPW
to within a few percent is not evidence for LPW at the level of a few percent;
it is compatible with a constant that is wrong by a factor of one and a half.
The pooled analysis does, however, establish a genuine bound: the universal
constant, if it exists, lies roughly in $[1.59,\,2.11]$, and no base deviates
from it detectably.

\section{The same error in a different guise}\label{sec:sameerror}

The Mersenne counting function invites a related error, which we present as a
worked example. Regressing $\pi_2(t)$ on $t$ at the event points, over the $50$
double-checked exponents (numbered $0,\dots,N-1$: the smallest exponent is the
origin, not an event), gives a slope of $2.6391$ and an intercept of $-2.9627$
with an OLS standard error of $0.4469$---apparently significant, and suggestive
of a secondary term in~\eqref{eq:LPW}. Nothing of the sort.

Two defects compound. The residuals of the cumulative count are strongly
autocorrelated, here $\rho_1=0.825$; the Newey--West covariance estimator with
eight lags raises the standard error to $0.7315$. More importantly, the
regression uses only the event points and thereby discards the empty interval
between the last exponent and the search frontier. The correct null model holds
the interval $[t_0,T]$ fixed, where $T=\ln X_{\mathrm{ver}}$, and conditions on
exactly $N-1=49$ events strictly to the right of $t_0$, which, by the
order-statistics property of the Poisson process given the number of events
(cf.\ the proof of Proposition~\ref{prop:unbiased}(b)), means $49$ independent
uniform points on $(t_0,T]$. Over $60\,000$ such replications the intercept has
median $-1.529$ and a $5$--$95\%$ range of $[-5.99,\,2.20]$, giving a two-sided
$p=0.570$ for the observed value. The secondary term is an artifact of
conditioning.

\subsection{Three consequences of one and the same conditioning}

\paragraph{(i) The bootstrap does not help.} The natural reaction to a
suspiciously small OLS standard error is to bootstrap over the event points. On
the base-2 data, a bootstrap over $20\,000$ resamples gives median $-2.948$,
standard deviation $0.385$ and a $95\%$ interval $[-3.809,\,-2.300]$ lying
entirely below zero; the proportion of negative values is $100\%$. This is not a
confirmation but a repetition of the same error: resampling the observed points
reproduces the variability within the realized configuration and cannot recover
the variability of the process that produced that configuration. The correct
null distribution has standard deviation $2.499$, that is, it is $6.5$ times
wider than the bootstrap one \emph{in terms of standard deviation} (and likewise
$6.5$ times in terms of the width of the $95\%$ interval).

\paragraph{(ii) Spurious underdispersion.} Comparing the observed
$\max_i|r_i|=6.347$, where $r_i=i-\hat\lambda(t_i-t_0)$, with a simulation of an
\emph{unconditioned} Poisson process suggests the conclusion that the counting
function is ``smoother than Poisson''. Against the correct null
distribution---conditioning on $N-1$ events and estimating $\lambda$ from the
last of them---the median is $5.117$ and the two-sided significance level is
$0.529$. The observation is entirely ordinary.

\paragraph{(iii) The ``mid-range deficit''.} Estimating $\lambda$ from the last
event forces the residual $r$ to vanish at both ends ($r(t_0)=r(t_N)=0$ under
the correct numbering), turning its trajectory into an analogue of a Brownian
bridge, for which a negative excursion in the middle of the range arises of its
own accord. The observed mean excursion over the middle half of the range is
$-3.13$ against a simulation median of $-0.012$ and $p=0.272$; the proportion of
replications with a negative mean excursion is $0.502$, so there is no
systematic deficit. The trajectory of the residual process itself, together with
the simulation band, is shown in Figure~\ref{fig:bridge}.

\begin{figure}[t]
\centering
\includegraphics[width=.78\textwidth]{fig_resid_bridge.pdf}
\caption{The residual process $r(t)$ for base $b=2$ together with the
$5$--$95\%$ quantile band of the simulations and the simulation median. The
trajectory stays inside the band.}
\label{fig:bridge}
\end{figure}

\subsection{Base 2 at a known frontier}

Base~2 is the only one whose frontier is published, and hence the only one for
which a scheme~B analysis is possible \emph{using the full exposure} rather than
only the conditional one. With $T=\ln X_{\mathrm{ver}}$ we obtain
$\hat\lambda=N/(T-t_0)=2.7960$ against $\eg/\ln2=2.5695$, that is,
$\hat\kappa_2=1.9381$ with an exact $95\%$ Garwood interval $[1.434,\,2.562]$
and a Poisson $p$-value of $0.593$ (doubled tail; under the mid-$p$ convention,
$0.545$---the convention must be stated here, since the difference is
appreciable). The conditional estimate on \emph{the same slice of data}---the
same $50$ double-checked exponents---gives $1.9046$; the difference is small
because the last exponent of that slice, $77\,232\,917$, happens to lie very
close to the frontier $81\,648\,221$. (For comparison, the conditional estimate
over all $52$ known exponents, reported in Table~\ref{tab:bases}, is $1.9215$;
it must not be compared with $\hat\kappa_2$, being a different sample.) This is
a coincidence of one realization, not a property of the method: the correction
of~\S\ref{sec:cost} concerns the average over the position of the last event
relative to the frontier.

We note in passing that the regime described here is exactly the one to which
we urge a transition for all bases: with $W_b$ known, conditioning on $D_b$
becomes unnecessary, and the precision increases by the very $3.7\%$ discussed
in~\S\ref{sec:exact}.

\subsection{The theoretical secondary term has the opposite sign}

There is an independent argument against a substantive reading of the intercept
$C_1$. If the secondary term in~\eqref{eq:LPW} originated in the Mertens
constant $\mathfrak M=0.26150\ldots$, its natural size would be
$(\eg/\ln2)\mathfrak M=+0.672$: the sign is positive, whereas what is observed is
$-2.963$. The nearest theoretical candidates not merely fail to reproduce the
observed magnitude---they predict the opposite sign. Independently of the
foregoing, this confirms that $C_1$ is a parameter of the fit and not of the
structure.

\section{Structural corrections}\label{sec:struct}

Consider the more general parametrization of the intensity
$\lambda_b=(\kappa/\ln b)\,C_b$ and four models for the factor:
\begin{align*}
\text{A:}&\quad C_b=1;\\
\text{B:}&\quad C_b=\prod_{p\mid b-1,\ p\ne2}\Bigl(\tfrac{p}{p-1}\Bigr)^{\alpha}
                 \prod_{p\mid b+1,\ p\ne2}\Bigl(\tfrac{p}{p-2}\Bigr)^{\beta};\\
\text{C:}&\quad C_b=\prod_{p\mid b^2-1}\Bigl(\tfrac{p}{p-1}\Bigr)^{\delta};\\
\text{D:}&\quad C_b=\exp\bigl(\theta_1\omega(b-1)+\theta_2\omega(b+1)+\theta_3\ln b\bigr).
\end{align*}
The exponent in model~C is written $\delta$ rather than $\gamma$ so as not to
confuse it with the Euler constant. The likelihood here is the \emph{conditional
Poisson} likelihood of Proposition~\ref{prop:cond},
$\ell_{\mathrm B}=\sum_b\bigl[K_b\ln\mu_b-\mu_b-\ln K_b!\bigr]$ with
$\mu_b=\kappa C_b D_b/\ln b$; the previous version used the scheme~A
likelihood~\eqref{eq:loglik}, which is incompatible with \S\ref{sec:scheme}. The
results are in Table~\ref{tab:struct}. No parametrization gives a significant
improvement; the BIC-based Bayes factors at $n=K=203$ overwhelmingly support the
base model.

\begin{table}[t]
\centering
\caption{Comparison of structural models under the conditional Poisson likelihood of
Proposition~\ref{prop:cond} (full form, including $-\ln K_b!$); LR and $p$ are
relative to model~A; BIC is computed with $n=K=203$.}
\label{tab:struct}
\small
\begin{tabular}{lrrrrrrr}
\toprule
Model & parameters & $\ln L$ & LR & $p$ & AIC & BIC & $P(\text{model}\mid\text{data})$\\
\midrule
A & 1 & $-35.753$ & --- & --- & 73.51 & 76.82 & 0.906\\
B & 3 & $-35.162$ & 1.18 & 0.554 & 76.32 & 86.26 & 0.008\\
C & 2 & $-35.458$ & 0.59 & 0.443 & 74.92 & 81.54 & 0.085\\
D & 4 & $-34.797$ & 1.91 & 0.591 & 77.59 & 90.85 & 0.001\\
\bottomrule
\end{tabular}
\end{table}

\begin{remark}[Sensitivity of the BIC]
The choice of effective sample size for the BIC is not obvious: the likelihood
has $B=15$ independent summands, but the information grows as $K=203$. With
$n=15$ the posterior probabilities are $0.661$ / $0.080$ / $0.229$ / $0.030$;
with $n=203$ they are $0.906$ / $0.008$ / $0.085$ / $0.001$. The conclusion
favors model~A under either choice, but with $n=203$ it is substantially
stronger. We report $n=K$ as matching the number of independent conditional
events.
\end{remark}

\paragraph{A fixed arithmetic correction.} We considered separately a model with
a \emph{fixed} (not estimated) factor
$C^{\mathrm{fix}}_b=\prod_{p\mid b-1,p\ne2}\frac{p}{p-1}\prod_{p\mid b+1,p\ne2}\frac{p}{p-2}$,
whose mean over the fifteen bases is $2.462$. It gives $\hat\kappa=0.759$. We
stress that comparing this value with $\eg$ is meaningless: in the
reparametrized model $\kappa$ has a different scale, and such a ``goodness-of-fit
test'' would only be testing whether the mean of $C_b$ equals one. Since models
A and $C^{\mathrm{fix}}$ are non-nested and have the same number of free
parameters, the correct comparison is directly by likelihood: $\ln L=-51.53$
against $-35.75$, a difference of $-15.78$. The fixed correction is decisively
rejected.

\paragraph{Logarithmic corrections.} A substantive correction is needed here.
The model $\lambda_b(t)=(\kappa/\ln b)(1+c_1/t+c_2/t^2)$ used in the previous
version is \emph{not an admissible intensity}: at the estimates obtained there,
$\hat c_1=+0.108$, $\hat c_2=-0.826$, the factor vanishes at $t=0.8564$ and is
negative to the left of that point, that is, for $n<2.355$; in particular, at
$t_0=\ln 2$ it equals $-0.5634$. The bases with $n_{\min}=2$
($b=2,6,10,12,18$) integrate the intensity over a region where it is negative,
and the measure $\int\lambda\,dt$ is negative on that stretch. That the
likelihood could be computed at all was a matter of chance: the smallest index
that is an event is $3$, that is, $t=1.0986>0.8564$, so no negative values fell
under the logarithm.

The correct parametrization of the same family is
\begin{equation}\label{eq:logcorr}
\lambda_b(t)=\frac{\kappa}{\ln b}\exp\Bigl(\frac{c_1}{t}+\frac{c_2}{t^2}\Bigr),
\end{equation}
which is positive identically and reduces to the null model at $c_1=c_2=0$. It
is estimated by the conditional point-process likelihood~\eqref{eq:condinhom}:
the sum $\sum\ln\lambda(t_i)$ is taken over the events \emph{excluding the last},
and the integral $\int\lambda$ runs from $t_0$ to the last event $t_0+D_b$, not
to the unknown frontier (Remark~\ref{rem:inhom}; the likelihood of the count
alone is inadmissible here: it discards the positions of the events, which are
the only thing that determines $c_1,c_2$). The result: $\hat\kappa=1.748$,
$\hat c_1=0.75$, $\hat c_2=-1.44$; the likelihood ratio against $c_1=c_2=0$ is
$2.03$ on two degrees of freedom, $p=0.36$ (tag~[S17]). There is no significant
improvement. Note, however, that $\hat c_1$ is close to the value $\ln2$
predicted by the heuristic itself; this topic is taken up in~\S\ref{sec:second}.

\section{The second-order term of the heuristic itself}\label{sec:second}

So far we have tested the constant-intensity model~\eqref{eq:lam}. But the LPW
heuristic itself predicts an intensity that is \emph{not} constant at
finite~$n$, and this changes the meaning of comparing $\hat\kappa$ with $\eg$.

\paragraph{Derivation.} Let $p$ be prime, $p\nmid b-1$, and let $q$ be a prime
divisor of $\Rb(p)$. Then $b$ has order $p$ modulo $q$, whence $q\equiv1\pmod p$,
and since $\Rb(p)$ is odd, $q\equiv1\pmod{2p}$. Hence $\Rb(p)$ has no prime
divisors below $2p+1$, and by Mertens' theorem the probability of primality is
increased by the factor $\prod_{q<2p}(1-1/q)^{-1}\approx\eg\ln(2p)$ compared with
a random number of the same size. This is the derivation
of~\cite{wagstaff1983,t5k}:
\begin{equation}\label{eq:probp}
\Pr\bigl(\Rb(p)\ \text{prime}\bigr)\approx\frac{\eg\ln(a_bp)}{p\ln b},\qquad a_b\ge2.
\end{equation}
Wagstaff's refinement for $b=2$ takes into account that $q\equiv\pm1\pmod8$ and
gives $a=2$ for $p\equiv3\pmod4$ and $a=6$ for $p\equiv1\pmod4$. For general $b$
the analogous necessary condition is that $b$ be a quadratic residue
modulo~$q$, since $b^p\equiv1$ and $p$ is odd; the smallest admissible candidate
$2kp+1$ determines $a_b=2k$. Multiplying~\eqref{eq:probp} by the density of
primes $1/\ln n$ and passing to $t=\ln n$, we get
\begin{equation}\label{eq:lam2}
\lambda_b(t)=\frac{\kappa}{\ln b}\Bigl(1+\frac{c_b}{t}\Bigr),\qquad
c_b=\langle\ln a_b\rangle,\qquad \kappa=\eg,
\end{equation}
and, for the counting function,
$\pi_b(x)\approx(\eg/\ln b)\,(\ln x+c_b\ln\ln x)$. The
asymptotics~\eqref{eq:LPW} are unaffected, but at $t\approx10$--$18$ the
correction $c_b/t$ amounts to $4$ to $17\%$---as much as or more than the
effects discussed in this paper. Averaging over primes $p\in[10^5,2\cdot10^5]$
gives $c_b$ from $1.23$ ($b=19$) to $1.74$ ($b=15$), and for $b=2$ the value
$1.242=\tfrac12(\ln2+\ln6)$, as it should be (script
\texttt{lpw\_second\_order.py}, tag~[S14]). The value $c=\ln2\approx0.693$ is a
lower bound that uses only $q\equiv1\pmod{2p}$. We stress that~\eqref{eq:probp}
is a heuristic; its second-order term is less reliable than the leading one, and
we treat $c$ as an external parameter in the range
$[\ln2,\ \langle\ln a_b\rangle]$.

\paragraph{Exact conditional inference for a given $c$.} By
Remark~\ref{rem:inhom}, conditionally on $D_b$ the interior events form a
Poisson process of intensity~\eqref{eq:lam2} on $(t_0,t_0+D_b)$. Since
$\int_{t_0}^{t_0+D_b}(1+c/t)\,dt=D_b+c\ln\bigl((t_0+D_b)/t_0\bigr)$,
Proposition~\ref{prop:cond} carries over verbatim:
\begin{equation}\label{eq:pivotc}
K\mid(D_b)_b\sim\Pois\bigl(\kappa S(c)\bigr),\qquad
S(c)=\sum_b\frac{D_b+c_b\ln\bigl((t_0+D_b)/t_0\bigr)}{\ln b},
\end{equation}
the estimator is $K/S(c)$, and the Garwood interval~\eqref{eq:garwood} applies
with $S$ replaced by $S(c)$.

\begin{table}[t]
\centering
\caption{Exact conditional inference for three values of the second-order
term~\eqref{eq:lam2}: the estimate $\hat\kappa=K/S(c)$ from~\eqref{eq:pivotc} and the
exact two-sided $p$ against $\kappa=\eg$. ``W'' is Wagstaff's
$c_b=\langle\ln a_b\rangle$. Truncation as in Table~\ref{tab:trunc}. Tag~[S16].}
\label{tab:second}
\small
\begin{tabular}{lrrrrrrr}
\toprule
truncation & $K$ & $\hat\kappa\,(c{=}0)$ & $\hat\kappa\,(c{=}\ln2)$ & $\hat\kappa$\,(W)
& $p\,(c{=}0)$ & $p\,(c{=}\ln2)$ & $p$\,(W)\\
\midrule
none & 203 & 1.838 & 1.625 & 1.453 & 0.672 & 0.200 & 0.003\\
$n>10$ & 177 & 1.871 & 1.712 & 1.578 & 0.531 & 0.630 & 0.111\\
$n>10^2$ & 141 & 1.830 & 1.704 & 1.595 & 0.771 & 0.636 & 0.201\\
$n>10^3$ & 110 & 1.823 & 1.715 & 1.620 & 0.832 & 0.738 & 0.344\\
$n>10^4$ & 75 & 1.789 & 1.695 & 1.612 & 1.000 & 0.723 & 0.421\\
\bottomrule
\end{tabular}
\end{table}

\paragraph{Results.} Table~\ref{tab:second} yields four conclusions.

\emph{(i) The second-order term is not determined by the data.} The profile MLE
of a common $c$ from the conditional likelihood~\eqref{eq:condinhom} on all the
data is $-0.40$ with a $95\%$ interval $[-0.69,\,0.44]$ (the left end is the
admissibility boundary $1+c/t_0>0$). But already at $n>10$ the interval widens to
$[-1.68,\,2.48]$, and at $n>10^3$ to $[-5.3,\,20]$ (tag~[S16]). The positions of
the events within the sequences carry almost no information about $c$: a shift
in $c$ is compensated by a shift in $\kappa$. The two-parameter
model~\eqref{eq:logcorr} gives no significant improvement either
(\S\ref{sec:struct}). The value of $c$ must therefore come from theory, not from
the data.

\emph{(ii) The inference about $\kappa$ depends on it more strongly than on the
observation scheme.} At $n>10$, moving from $c=0$ to $c=\ln2$ lowers the estimate
by $8.5\%$, and moving to $c=\langle\ln a_b\rangle$ lowers it by $15.7\%$. This
exceeds the price of the choice of observation scheme ($6.9\%$,
\S\ref{sec:cost}) and is two to three times the residual gap between $(M-B)/S$
and $\eg$ ($3.2\%$). Comparing the constant-intensity estimate with $\eg$ is not
a test of LPW at finite $n$: LPW itself predicts that such an estimate lies
\emph{above} $\eg$.

\emph{(iii) LPW with its own second-order term is consistent with the data.} For
$c=\ln2$ and $n>10$ the exact $p$ lies within $0.63$--$0.74$. For
$c=\langle\ln a_b\rangle$ it is $0.11$--$0.42$. On the full sample with
$c=\langle\ln a_b\rangle$ one obtains $p=0.003$. But the full sample includes
$n\le10$, where expansion~\eqref{eq:lam2} is meaningless: at $n=2$ the factor
$1+c/t$ equals $2$ even for $c=\ln2$, and the indices $2,3,5,7$ themselves are
dictated by the arithmetic of small numbers. We do not regard this result as
evidence against LPW; it is a warning against using small indices when $c\ne0$.

\emph{(iv) The trend under truncation is explained by the second-order term.}
For $c=0$ the estimate decreases monotonically as the lower bound grows, from
$1.871$ at $n>10$ to $1.789$ at $n>10^4$ ($-4.4\%$). This is exactly what one
expects from~\eqref{eq:lam2}: the mean of $1/t$ over the exposure falls as $t_0$
grows. For $c=\ln2$ the same four estimates lie within $1.695$--$1.715$ (a range
of $1.2\%$). The subsamples are nested, so this observation is illustrative
rather than a separate test.

\paragraph{Consequence for the folklore.} The term $c_b\ln\ln x$ in the counting
function also explains the observation with which this work began: the
normalized frequencies lie below $\emg$ and converge to it slowly from below.
Under~\eqref{eq:lam2} the slope of $t_k$ against $k$ is
$\emg\ln b/(1+c_b\langle1/t\rangle)$, so the normalized frequency is depressed by
the factor $1+c_b\langle1/t\rangle$ and tends to $\emg$ as $1/t$, that is, as
$1/\ln n$---slowly indeed. In~\S\ref{sec:notbias} we showed that fitting up to
the right endpoint does not create the effect; here we see where it comes from:
from the heuristic itself, not from~$\kappa\ne\eg$.

\paragraph{Summary.} The question ``are the data compatible with LPW?'' splits
into two: one about the leading term, $\kappa=\eg$, and one about the
second-order term, $c_b$. The latter is not determined by the data, but it
affects the answer to the former more strongly than everything analysed in
\S\ref{sec:scheme}. For any $c\in[\ln2,\,\langle\ln a_b\rangle]$ and $n>10$, LPW
is not rejected. Refining the second-order term is a problem for number theory,
not for statistics. Without it, a quantitative test of LPW more precise than
$\pm10\%$ remains out of reach even as the frontiers grow: the term $c_b/t$
decays only as $1/\ln n$.

\section{Conclusions}

\begin{enumerate}
\item The estimators of the normalized frequency of generalized repunit primes
that are in common use are \emph{not} biased by the fitting scheme: both the
endpoint form and the OLS slope are exactly unbiased under scheme~A and unbiased
up to an error exponentially small in $\mu$ under scheme~B (at most $0.73\%$ for
the endpoint form and $2.3\%$ for OLS when $\mu\ge7$). Only the reciprocal is
biased, and that is a consequence of convexity. The common explanation of ``slow
convergence from below'' in terms of censoring fails under both schemes.

\item The genuine defect of current practice is that different forms of the
estimator disagree by up to $29.5\%$, and that point values are published
without confidence intervals even though the OLS standard errors in this problem
are understated: the Newey--West correction increases them by a factor of $1.6$,
and the correct conditional null distribution is $5.6$ times wider than the
nominal one (\S\ref{sec:sameerror}). The classification of bases into ``lucky''
and ``unlucky'' flips when the form of the estimator is changed on unchanged
data; for two of the fifteen bases even the sign of the deviation from $\emg$
flips.

\item The main result is the exact conditional likelihood for scheme~B. The
unrecorded frontiers are Neyman--Scott nuisance parameters, and the profile
estimator $M/S$ is inconsistent. The position $D_b$ of the last event is
sufficient for the frontier, and conditionally on it $N_b-1$ is distributed as
$\Pois(\lambda_bD_b)$ under any rule for choosing the frontier that does not
depend on the finds. The same result is obtained by integrating over the
frontier against the invariant measure $dW_b$. Hence
$M-B\mid D\sim\Pois(\kappa S)$: the correct pivot under a fixed frontier is
Poisson rather than $\Gam(M,1)$, the correct pooled estimator is $(M-B)/S$
rather than $(M-1)/S$, and confidence intervals are available in closed form
without simulation-based inversion. The conditional likelihood costs one event
per base, not one event in total---an exact analogue of the loss of one degree
of freedom per group in the Neyman--Scott example. Distinguishing such schemes
goes back to reliability theory; what is new here is the case of an
undocumented frontier, where the exposure cannot be taken from a protocol and
has to be reconstructed from the finds (Remark~\ref{rem:priority}).

\item This gives $\kappa=1.838$ with an exact $95\%$ interval
$[1.594,\,2.109]$ and $p=0.672$ against $\eg$. The reported ten-percent excess
over LPW---including our own in an earlier version of this work---is an artifact
of the observation scheme. Under scheme~B the old Gamma pivot rejects a true
hypothesis in $11$--$17\%$ of cases. Moving to scheme~B also independently
removes the systematic trend of the pooled estimate under truncation in
$n_{\min}$.

\item The inference rests on a single assumption (A1): the frontier was chosen
without regard to the finds. It cannot be checked from the data: $(M-1)/S$ and
$(M-B)/S$ are functions of the same statistics and differ by $6.9\%$. Without
(A1) only the upper bound survives rigorously (Wald's identity); the rule ``stop
after a long stretch without finds'' biases even $(M-B)/S$ upwards, so there is
no two-sided bracket. None of the violations of (A1) analysed here creates an
excess over LPW. The four documented frontiers point, if anywhere, to a
deviation towards stopping close to a find rather than ``after a drought''; the
hybrid estimate with the known frontiers for $b=2$ and $b=10$ is $1.853$. The
practical consequence is the need to publish the search frontiers $L_b$ (and,
with them, the protocols of sieving completeness), as GIMPS does for $b=2$.
Until that is done, any pooled analysis of these sequences rests on an
unverifiable assumption whose price is at least $7\%$, that is, twice the
residual gap with LPW.

\item The LPW conjecture is not rejected by the available data, and
base-dependent corrections are not required. But both of these are statements
about the low sensitivity of the tests: the conditional multinomial homogeneity
test resolves only a spread of $\kappa_b$ of order $\pm37\%$, and the pooled
test only deviations of the constant of at least $21.3\%$ (two-sided variant).
Attaining a sensitivity of $10\%$ would require the search frontiers to grow by
a factor of about $5\cdot10^{19}$.

\item The LPW heuristic itself predicts at finite $n$ an intensity
$(\kappa/\ln b)(1+c_b/t)$ with $c_b\in[\ln2,\,1.74]$, so comparing the
constant-intensity estimate with $\eg$ is not a test of LPW. The second-order
term is not determined by the data but shifts the estimate of $\kappa$ by
$8$--$16\%$; once it is taken into account, LPW is consistent with the data, the
trend under truncation disappears, and the ``slow convergence from below'' is
explained. Refining $c_b$ is a problem for number theory; until it is solved, a
test of LPW more precise than $\pm10\%$ is out of reach.
\end{enumerate}

\appendix
\section{Verification of the source data}\label{app:data}

All twenty sequences (fifteen of the main set and five from the check of the
selection limit) were cross-checked against the OEIS entries on 27 August 2026.
Since only the triples $(N_b,n_{\min},n_{\max})$ enter the pooled analysis,
Table~\ref{tab:verify} gives, for the main set, the number of known primes---from
which $N_b$ is obtained by subtracting one---and both extreme indices; the full
lists of indices are contained in the script \texttt{analysis\_v2.py}.

\begin{table}[h]
\centering
\caption{Cross-check against OEIS. The column ``last addition'' gives the
provenance of the largest known index.}
\label{tab:verify}
\small
\begin{tabular}{rlrrrl}
\toprule
$b$ & OEIS & primes & $n_{\min}$ & $n_{\max}$ & last addition\\
\midrule
2 & A000043 & 52 & 2 & 136\,279\,841 & GIMPS, Oct.\ 2024\\
3 & A028491 & 23 & 3 & 8\,530\,117 & P.~Bourdelais\\
5 & A004061 & 19 & 3 & 3\,300\,593 & P.~Bourdelais\\
6 & A004062 & 17 & 2 & 3\,360\,347 & P.~Bourdelais\\
7 & A004063 & 10 & 5 & 1\,264\,699 & P.~Bourdelais\\
10 & A004023 & 11 & 2 & 8\,177\,207 & Propper and Batalov, May 2021\\
11 & A005808 & 13 & 17 & 1\,868\,983 & P.~Bourdelais\\
12 & A004064 & 14 & 2 & 769\,543 & P.~Bourdelais, Dec.\ 2014\\
13 & A016054 & 13 & 5 & 1\,503\,503 & P.~Bourdelais\\
14 & A006032 & 11 & 3 & 1\,724\,417 & P.~Bourdelais, 13 Jul.\ 2026\\
15 & A006033 & 10 & 3 & 639\,833 & P.~Bourdelais, 22 Apr.\ 2019\\
17 & A006034 & 12 & 3 & 1\,990\,523 & P.~Bourdelais, 3 Aug.\ 2020\\
18 & A133857 & 8 & 2 & 1\,270\,141 & P.~Bourdelais\\
19 & A006035 & 11 & 19 & 209\,359 & P.~Bourdelais, 27 Aug.\ 2010\\
20 & A127995 & 9 & 3 & 984\,349 & P.~Bourdelais\\
\bottomrule
\end{tabular}
\end{table}

A separate warning to the reader who checks the data against secondary sources:
summary tables (the ``Repunit primes'' page in the OEIS Wiki, the article ``List
of repunit primes'' in Wikipedia) are updated with a delay and, at the time of
writing, lagged the OEIS entries themselves by one term for the bases
$b=14,15,17,18,20$. The OEIS entries should be treated as authoritative.

\subsection*{Independent recomputation of the prefixes}

Cross-checking against OEIS confirms that the listed indices give primes, but it
does not guarantee \emph{completeness}: a missing term would shift $N_b$, and
with it both pooled estimates. We closed this gap by direct recomputation. For
each of the twenty bases $2\le b\le26$ that are not perfect powers, an
independent implementation\footnote{A hybrid CPU/GPU searcher: segmented sieving
by small divisors, trial factoring in Montgomery arithmetic, PRP testing via GMP
and GWNUM/IBDWT. The source code and the run logs are in the same repository,
scripts \texttt{verify\_sequences.sh} and \texttt{compare\_verify.py}.} tested
\emph{all} indices $n<10^4$ consecutively, not only those listed in OEIS. Every
negative verdict was recomputed by a second, independent backend (a blanket
double-check, as in GIMPS) so as to rule out a lost find.

The result: $153$ terms confirmed, none spurious, none missing, no disagreement
between backends---across all twenty bases. Coverage is uneven across bases
(from one term at $b=18$ to twenty-two at $b=2$): the main set $b\le20$ accounts
for $128$ confirmed indices out of $233$ known. Above $10^4$ a complete
recomputation is computationally out of reach, and there we still rely on OEIS.

We note that this recomputation was not a formality: on the first run the
independent implementation declared $R_{3181}(23)$ composite, disagreeing with
A204940. Analysis showed an error in the implementation itself---a spike of FFT
rounding error on the last iteration, after the last periodic check---and not in
the OEIS data. After the fix no disagreements remained. We report this episode as
an illustration of the thesis of \S\ref{sec:scheme}: a false ``composite'' is
invisible without recomputation, and that is exactly why frontiers and
verification protocols should be published.

\section*{Declarations}

\subsection*{Acknowledgments}
The author thanks the Great Internet Mersenne Prime Search (GIMPS) project for
providing the Mersenne prime data through the open PrimeNet report
(\url{https://www.mersenne.org/}); the Mersenne exponents were retrieved on
27 August 2026. The author is grateful to the OEIS contributors who have
maintained the generalized repunit prime sequences for many years, and to Jaume
Gotanegra (MathOverflow, \url{https://mathoverflow.net/questions/514719}) for
pointing out that the conditional likelihood~\eqref{eq:condlaw} is also obtained
by integrating out the frontier against $dW_b$ (Remark~\ref{rem:integrated}).

\subsection*{Funding}
This research received no external funding.

\subsection*{Competing interests}
The author declares no competing interests.

\subsection*{Data and code availability}
All sequences used are listed in Table~\ref{tab:data}.

\emph{Data.} The script \texttt{fetch\_data.sh} downloads the OEIS b-files for
all fifteen sequences of the main set, and \texttt{fetch\_extra.sh} does the
same for the five bases $21\le b\le26$ used in the check of the selection limit
(\S\ref{sec:bases}). The downloaded files are included in the repository, so the
analysis reproduces without network access.

\emph{Analysis.} The script \texttt{analysis\_v2.py} (numpy, scipy, matplotlib);
every number in this paper is tagged in its output (\texttt{results\_v2.txt})
with a tag of the form~[S3] identifying the corresponding section or table. The
pooled analysis, all confidence intervals, the homogeneity test, the power
calculations and the structural models depend only on the triples
$(N_b,n_{\min},n_{\max})$; the full lists of indices are needed only for the
$\hat G^{\mathrm{OLS}}_b$ column of Table~\ref{tab:bases}, for
\S\ref{sec:sameerror}, for the truncation and for the variable-intensity models.
The script reads them from the b-files in the \texttt{data/} directory (the
loader is shared with \texttt{analysis.py} of version~1). The random number
generator is initialized with a fixed seed ($20260914$), so all Monte Carlo
quantities reproduce bit for bit. The script also produces all four figures (directory \texttt{figs\_v2/}, with
English labels in \texttt{figs\_v2/en/}).
Table~\ref{tab:stop} (stopping rules that violate (A1)) is produced by a separate
script, \texttt{stopping\_rules.py} (tag~[S13], seed $20260924$). The
second-order term (\S\ref{sec:second}), model~\eqref{eq:logcorr} and the
homogeneity test at $B=20$ are computed by \texttt{lpw\_second\_order.py}
(tags~[S14]--[S17]), and the check of (A1) against the documented frontiers and
the hybrid estimate by \texttt{frontiers.py} (tag~[S18]; the sources of the
frontiers are listed in the script itself). The output of all three scripts is
saved in \texttt{results\_v3.txt}.

\emph{Independent recomputation of the sequences.}
\texttt{verify\_sequences.sh} runs the searcher over all indices $n<10^4$ for
each base, and \texttt{compare\_verify.py} compares the result with the OEIS
b-files, returning a non-zero exit code on any disagreement
(Appendix~\ref{app:data}).

\emph{A caveat on the Monte Carlo quantities.} The main result~\eqref{eq:resB},
the intervals of Table~\ref{tab:bases}, Table~\ref{tab:trunc} and
Table~\ref{tab:power} are \emph{deterministic}: they are obtained by exact
inversion of the Poisson distribution and do not depend on the random number
generator. Only Tables~\ref{tab:mc}, \ref{tab:calib} and~\ref{tab:stop}, the
exact $p$-value of the homogeneity test and its power curve, and the null
distributions of \S\ref{sec:sameerror} depend on simulation; their accuracy is
stated where they appear. This is an essential difference from the previous
version, where the main interval estimate itself was simulation-based.

Repository: \url{https://github.com/Wilps93/repunit-hunt}, directory
\texttt{paper/}. The file tree and the order of execution are described in
\texttt{paper/README.md}.

\begin{thebibliography}{99}

\bibitem{wagstaff1983} S.~S. Wagstaff, Jr., \emph{Divisors of Mersenne
numbers}, Math. Comp. \textbf{40} (1983), no.~161, 385--397.
\href{https://doi.org/10.1090/S0025-5718-1983-0679454-X}{DOI:
10.1090/S0025-5718-1983-0679454-X}.

\bibitem{pomerance1981} C.~Pomerance, \emph{Recent developments in primality
testing}, Math. Intelligencer \textbf{3} (1981), no.~3, 97--105.
\href{https://doi.org/10.1007/BF03022861}{DOI: 10.1007/BF03022861}.

\bibitem{t5k} \emph{Heuristics: deriving the Wagstaff Mersenne conjecture}, The
PrimePages: prime number research and records.
\url{https://t5k.org/mersenne/heuristic.html} (accessed 27 August 2026).

\bibitem{dubner1993} H.~Dubner, \emph{Generalized repunit primes}, Math. Comp.
\textbf{61} (1993), no.~204, 927--930.
\href{https://doi.org/10.1090/S0025-5718-1993-1185243-9}{DOI:
10.1090/S0025-5718-1993-1185243-9}.

\bibitem{bourdelais2009} P.~Bourdelais, \emph{A generalized repunit
conjecture}, NMBRTHRY mailing list, 25 June 2009.
\url{https://listserv.nodak.edu/cgi-bin/wa.exe?A2=NMBRTHRY;417ab0d6.0906}
(accessed 25 August 2026).

\bibitem{oeis} OEIS Foundation Inc., \emph{The On-Line Encyclopedia of Integer
Sequences}. \url{https://oeis.org} (accessed 27 August 2026).

\bibitem{wagstaff2025} S.~S. Wagstaff, Jr., \emph{Two Mersenne prime
conjectures}, J. Integer Seq. \textbf{28} (2025), article 25.7.2.
\url{https://cs.uwaterloo.ca/journals/JIS/VOL28/Wagstaff/wag14.html}
(accessed 24 September 2026).

\bibitem{dominguez2026} J.~Dominguez, \emph{Arithmetic bias in the distribution
of Mersenne prime exponents and the divisor structure of $p-1$}, preprint,
arXiv:2603.08994 (2026). \url{https://arxiv.org/abs/2603.08994}
(accessed 27 August 2026).

\bibitem{neymanscott1948} J.~Neyman and E.~L. Scott, \emph{Consistent estimates
based on partially consistent observations}, Econometrica \textbf{16} (1948),
no.~1, 1--32.
\href{https://doi.org/10.2307/1914288}{DOI: 10.2307/1914288}.

\bibitem{andersen1970} E.~B. Andersen, \emph{Asymptotic properties of
conditional maximum-likelihood estimators}, J. Roy. Statist. Soc. Ser.~B
\textbf{32} (1970), no.~2, 283--301.
\href{https://doi.org/10.1111/j.2517-6161.1970.tb00842.x}{DOI:
10.1111/j.2517-6161.1970.tb00842.x}.

\bibitem{jha2026} A.~Jha, \emph{The Poisson tail conjecture for primes in short
intervals}, preprint, arXiv:2605.23014 (2026).
\url{https://arxiv.org/abs/2605.23014} (accessed 27 August 2026).

\bibitem{gimps} Great Internet Mersenne Prime Search (GIMPS), \emph{PrimeNet};
Prime95/MPrime software, version 30.8. \url{https://www.mersenne.org/}
(accessed 27 August 2026; the \emph{Milestones} report was read at
27 August 2026, 04:30~UTC).

\bibitem{epstein1954trunc} B.~Epstein, \emph{Truncated life tests in the
exponential case}, Ann. Math. Statist. \textbf{25} (1954), no.~3, 555--564.
\href{https://doi.org/10.1214/aoms/1177728723}{DOI: 10.1214/aoms/1177728723}.

\bibitem{epsteinsobel1954} B.~Epstein and M.~Sobel, \emph{Some theorems
relevant to life testing from an exponential distribution}, Ann. Math. Statist.
\textbf{25} (1954), no.~2, 373--381.
\href{https://doi.org/10.1214/aoms/1177728793}{DOI: 10.1214/aoms/1177728793}.

\bibitem{bartholomew1963} D.~J. Bartholomew, \emph{The sampling distribution of
an estimate arising in life testing}, Technometrics \textbf{5} (1963), no.~3,
361--374.
\href{https://doi.org/10.1080/00401706.1963.10490104}{DOI:
10.1080/00401706.1963.10490104}.

\bibitem{sengupta2009} S.~N. Sengupta, B.~K. Sinha, and S.~B. Bagchi,
\emph{Estimation of mean life and reliability for exponential life distribution
using time censored sample data}, Statist. Probab. Lett. \textbf{79} (2009),
no.~7, 899--906.
\href{https://doi.org/10.1016/j.spl.2008.11.026}{DOI:
10.1016/j.spl.2008.11.026}.

\bibitem{kivela2015} M.~Kivel\"{a} and M.~A. Porter, \emph{Estimating interevent
time distributions from finite observation periods in communication networks},
Phys. Rev. E \textbf{92} (2015), no.~5, 052813.
\href{https://doi.org/10.1103/PhysRevE.92.052813}{DOI:
10.1103/PhysRevE.92.052813}.

\bibitem{straussSadler1989} D.~Strauss and P.~M. Sadler, \emph{Classical
confidence intervals and Bayesian probability estimates for ends of local taxon
ranges}, Math. Geol. \textbf{21} (1989), no.~4, 411--427.
\href{https://doi.org/10.1007/BF00897326}{DOI: 10.1007/BF00897326}.

\bibitem{marshall1990} C.~R. Marshall, \emph{Confidence intervals on
stratigraphic ranges}, Paleobiology \textbf{16} (1990), no.~1, 1--10.
\href{https://doi.org/10.1017/S0094837300009672}{DOI:
10.1017/S0094837300009672}.

\bibitem{solow1993} A.~R. Solow, \emph{Inferring extinction from sighting data},
Ecology \textbf{74} (1993), no.~3, 962--964.
\href{https://doi.org/10.2307/1940821}{DOI: 10.2307/1940821}.

\bibitem{robsonwhitlock1964} D.~S. Robson and J.~H. Whitlock, \emph{Estimation
of a truncation point}, Biometrika \textbf{51} (1964), no.~1--2, 33--39.
\href{https://doi.org/10.1093/biomet/51.1-2.33}{DOI:
10.1093/biomet/51.1-2.33}.

\bibitem{liu2023} Z.~Liu, U.~Bhandaram, and N.~Garg, \emph{Quantifying spatial
under-reporting disparities in resident crowdsourcing}, Nat. Comput. Sci.
\textbf{4} (2024), no.~1, 57--65.
\href{https://doi.org/10.1038/s43588-023-00572-6}{DOI:
10.1038/s43588-023-00572-6}.

\bibitem{lancaster2000} T.~Lancaster, \emph{The incidental parameter problem
since 1948}, J. Econometrics \textbf{95} (2000), no.~2, 391--413.
\href{https://doi.org/10.1016/S0304-4076(99)00044-5}{DOI:
10.1016/S0304-4076(99)00044-5}.

\bibitem{kalbfleischsprott1970} J.~D. Kalbfleisch and D.~A. Sprott,
\emph{Application of likelihood methods to models involving large numbers of
parameters} (with discussion), J. Roy. Statist. Soc. Ser.~B \textbf{32} (1970),
no.~2, 175--208.
\href{https://doi.org/10.1111/j.2517-6161.1970.tb00830.x}{DOI:
10.1111/j.2517-6161.1970.tb00830.x}.

\bibitem{berger1999} J.~O. Berger, B.~Liseo, and R.~L. Wolpert,
\emph{Integrated likelihood methods for eliminating nuisance parameters},
Statist. Sci. \textbf{14} (1999), no.~1, 1--28.
\href{https://doi.org/10.1214/ss/1009211804}{DOI: 10.1214/ss/1009211804}.

\bibitem{lindsay1980} B.~G. Lindsay, \emph{Nuisance parameters, mixture models,
and the efficiency of partial likelihood estimators}, Philos. Trans. Roy. Soc.
London Ser.~A \textbf{296} (1980), no.~1427, 639--662.
\href{https://doi.org/10.1098/rsta.1980.0197}{DOI: 10.1098/rsta.1980.0197}.

\bibitem{garwood1936} F.~Garwood, \emph{Fiducial limits for the Poisson
distribution}, Biometrika \textbf{28} (1936), no.~3--4, 437--442.
\href{https://doi.org/10.1093/biomet/28.3-4.437}{DOI:
10.1093/biomet/28.3-4.437}.

\end{thebibliography}

\end{document}
